% Options for packages loaded elsewhere
\PassOptionsToPackage{unicode}{hyperref}
\PassOptionsToPackage{hyphens}{url}
\documentclass[
]{article}
\usepackage{xcolor}
\usepackage{amsmath,amssymb}
\setcounter{secnumdepth}{5}
\usepackage{iftex}
\ifPDFTeX
  \usepackage[T1]{fontenc}
  \usepackage[utf8]{inputenc}
  \usepackage{textcomp} % provide euro and other symbols
\else % if luatex or xetex
  \usepackage{unicode-math} % this also loads fontspec
  \defaultfontfeatures{Scale=MatchLowercase}
  \defaultfontfeatures[\rmfamily]{Ligatures=TeX,Scale=1}
\fi
\usepackage{lmodern}
\ifPDFTeX\else
  % xetex/luatex font selection
\fi
% Use upquote if available, for straight quotes in verbatim environments
\IfFileExists{upquote.sty}{\usepackage{upquote}}{}
\IfFileExists{microtype.sty}{% use microtype if available
  \usepackage[]{microtype}
  \UseMicrotypeSet[protrusion]{basicmath} % disable protrusion for tt fonts
}{}
\makeatletter
\@ifundefined{KOMAClassName}{% if non-KOMA class
  \IfFileExists{parskip.sty}{%
    \usepackage{parskip}
  }{% else
    \setlength{\parindent}{0pt}
    \setlength{\parskip}{6pt plus 2pt minus 1pt}}
}{% if KOMA class
  \KOMAoptions{parskip=half}}
\makeatother
\usepackage{longtable,booktabs,array}
\usepackage[margin=1in]{geometry}
\allowdisplaybreaks
\setlength{\emergencystretch}{3em}
\usepackage{caption}
\captionsetup[table]{skip=6pt}
\newcounter{none} % for unnumbered tables
\usepackage{calc} % for calculating minipage widths
% Correct order of tables after \paragraph or \subparagraph
\usepackage{etoolbox}
\makeatletter
\patchcmd\longtable{\par}{\if@noskipsec\mbox{}\fi\par}{}{}
\makeatother
\setlength{\emergencystretch}{3em} % prevent overfull lines
\providecommand{\tightlist}{%
  \setlength{\itemsep}{0pt}\setlength{\parskip}{0pt}}
\usepackage{bookmark}
\IfFileExists{xurl.sty}{\usepackage{xurl}}{} % add URL line breaks if available
\urlstyle{same}
% fallback for those not using the hyperref driver hyperxmp:
\makeatletter
\@ifundefined{xmpquote}{\newcommand{\xmpquote}[1]{#1}}{}
\makeatother
\hypersetup{
  pdftitle={Related prior work for Expansion},
  pdfauthor={Wren, an Open Math OS Agent},
  hidelinks,
  pdfcreator={LaTeX via pandoc}}

\title{Related prior work for Expansion}
\author{Wren, an Open Math OS Agent}
\date{Compiled \today}

\begin{document}
\maketitle
\tableofcontents
\newpage

\textbf{Scope.} I read the draft at \texttt{dpthurst/expansion}, commit
\texttt{73fe2f1} (2026-10-18), including \texttt{alternatives.tex} and
the unintegrated maximum-growth material in \texttt{clippings.tex}. This
memo is a literature map for the two threads you named. Every work it
cites is collected once, by thread, in the bibliography at the end
(Section 9). One distinction runs through the whole report and is easy
to lose: \textbf{topological} invariants (invariant under topological
conjugacy --- what your paper produces), \textbf{metric-bound}
invariants (invariant only under bi-Lipschitz changes of the metric),
and \textbf{smooth} invariants (which need a differentiable structure).
Section 2.1 sets out the three kinds and places the literature's
quantities in them; each later subsection says which kind it is
discussing. It deliberately does not address any open question in the
draft, and it does not try to prove or disprove anything in it.

\textbf{One-paragraph summary.} Both threads already have substantial
literatures, and the two are more connected than the draft's separation
suggests. As far as I found, the closest published invariant to your
\(h_{\min}\) is the exponential decay rate of \textbf{Lebesgue numbers}
of the refined covers (Sun 2012): your in-radius
\(r(\mathcal U;\mathcal U^{(n)})\) is essentially the reciprocal of the
Lebesgue number \(\delta(\mathcal U^{(n)})\), so, up to the comparison
between in-radius and Lebesgue number, \(h_{\min}\) is minus Sun's
\(h_L\). The closest smooth invariant to your \(\lambda_{\exp}\) is
\textbf{Hirsch's logarithmic expansion rate}
\(E(f,K)=\sup_n \frac1n\log EC(f^n,K)\) and \textbf{Schreiber's theorem}
that, for \(C^1\) maps, \(E\) is the infimum of the smallest Lyapunov
exponent over ergodic invariant measures. The closest circle of ideas to
your visual-metric theorem is the \textbf{entropy--dimension--expansion
inequality} (Fathi 1989; Dai--Zhou--Geng 1998; Misiurewicz 2004;
Fujita--Kato--Matsumoto 2010; Roth--Roth 2020). In its sharp form it
says that for a metric in which \(f\) expands small distances by exactly
\(\lambda\), \[
\dim X\cdot\log\lambda \;\le\; D_d(X)\log\lambda \;\le\; h_{\mathrm{top}}(f) \;\le\; D_d(X)\log\operatorname{Lip}_d(f),
\] where \(D_d(X)\) is the box dimension in the metric \(d\). Your
\(h_{\min}\) is the topological minimal distance growth,
\(h_{\mathrm{top}}\) is the volume growth, and
\(\lambda_{\exp}=e^{h_{\min}}\) is the maximal expanding constant. The
draft's theorem \(\log\lambda_{\exp}=h_{\min}\) is the topological
counterpart of Schreiber's smooth theorem. At the other end, the closest
published objects are Kifer's pointwise maximum and minimum
characteristic exponents in metric spaces (1983) --- the pointwise
analogue of your \(h_{\max}\), but depending on the metric --- and the
infimum of the Lipschitz constant in a conjugacy class of interval maps
(Bobok--Roth 2019), the closest relative of \(\lambda_{\mathrm{Lip}}\);
Section 8 maps that side, and nothing I found asserts or denies
\(h_{\max}=\log\lambda_{\mathrm{Lip}}\).

\begin{center}\rule{0.5\linewidth}{0.5pt}\end{center}

\section{The six papers I would read
first}\label{the-six-papers-i-would-read-first}

\begin{enumerate}
\def\labelenumi{\arabic{enumi}.}
\item
  \textbf{Peng Sun, ``Exponential decay of Lebesgue numbers,''}
  \emph{Discrete Contin. Dyn. Syst.} \textbf{32} (2012), 3773--3785;
  arXiv:1002.0549. Defines \[
  h_L^\pm(f,\mathcal U)=\limsup/\liminf_{n\to\infty}-\frac1n\log\delta\bigl(\textstyle\bigvee_{k<n}f^{-k}\mathcal U\bigr),
  \] where \(\delta\) is the Lebesgue number, and
  \(h_L^\pm(f)=\sup_{\mathcal U}h_L^\pm(f,\mathcal U)\). Proves
  \(\dim_B(X)\,h_L^-(f)\ge h(f)\) and \(\dim_H(X)\,h_L^+(f)\ge h(f)\);
  for Lipschitz \(f\), \(h_L^+\le\log L(f)\); \(h_L\) is bi-Lipschitz
  invariant but not a topological invariant. This is the closest
  published invariant to your \(h_{\min}\).
\item
  \textbf{Sebastian J. Schreiber, ``Expansion rates and Lyapunov
  exponents,''} \emph{Discrete Contin. Dyn. Syst.} \textbf{3} (1997),
  433--438; DOI 10.3934/dcds.1997.3.433. For a \(C^1\) endomorphism
  \(f\) of a Riemannian manifold and a compact positively invariant
  \(K\) on which \(f\) is immersive, defines the expansion constant
  \(EC(f,K)\) (the largest \(\rho\ge0\) such that for every
  \(0<\rho^*<\rho\) there is \(\delta>0\) with
  \(d(fx,fy)\ge\rho^*d(x,y)\) whenever \(x\in K\), \(y\in M\) and
  \(d(x,y)\le\delta\)) and Hirsch's logarithmic expansion rate
  \(E(f,K)=\sup_n\frac1n\log EC(f^n,K)\). Proves \[
  E(f,K)=\inf_\mu E(f,\mu),
  \] the infimum of the smallest Lyapunov exponent over ergodic
  invariant measures supported on \(K\). This is the smooth analogue of
  \(\log\lambda_{\exp}=h_{\min}\).
\item
  \textbf{Chikara Fujita, Hisao Kato, Masahiro Matsumoto, ``Fractal
  metrics of Ruelle expanding maps and expanding ratios,''}
  \emph{Topology Appl.} \textbf{157} (2010), 615--628; DOI
  10.1016/j.topol.2009.11.002. For a Ruelle expanding map \(f\) of a
  compactum \(X\) and every \(s>1\), constructs an admissible metric
  \(d\) and an \textbf{exact} expansion ratio \(\lambda\in(1,s)\):
  \(d(fx,fy)=\lambda d(x,y)\) for \(d(x,y)\le\epsilon\). In that metric,
  \[
  \dim_H(X,d)=\underline D_d(X)=D_d(X)=\frac{h(f)}{\log\lambda}.
  \] Proposition 4.1 gives, for \(1<\lambda_2\le\lambda_1\) and
  \(\lambda_2 d(x,y)\le d(fx,fy)\le\lambda_1 d(x,y)\) for small
  \(d(x,y)\), \[
  D_d(X)\log\lambda_2\le h(f)\le D_d(X)\log\lambda_1.
  \] Remark 4.4 deduces \(\log\lambda\le h(f)/\dim X\) for a strictly
  expanding ratio \(\lambda\); hence the set of expanding ratios is
  bounded above by \(e^{h/\dim X}\). This is the closest ``set of
  possible expanding constants'' result.
\item
  \textbf{Albert Fathi, ``Expansiveness, hyperbolicity and Hausdorff
  dimension,''} \emph{Comm. Math. Phys.} \textbf{126} (1989), 249--262;
  DOI 10.1007/BF02125125. For a compact hyperbolic set \(K\) of a
  \(C^1\) diffeomorphism, with uniform contraction rate \(\rho<1\) for
  \(f\) on the stable manifolds and for \(f^{-1}\) on the unstable
  manifolds, \[
  \mathrm{HD}(K)\le \frac{2\,\mathrm{ent}(f|_K)}{\log(1/\rho)}.
  \] The same argument with no stable direction gives, for an expanding
  map, \(\mathrm{HD}(K)\le h(f)/\log\lambda\). The paper also constructs
  a ``hyperbolic metric'' and uses the inequality to prove
  Ma\~{n}\'{e}'s theorem that expansive homeomorphisms have finite
  topological dimension.
\item
  \textbf{Samuel Roth, Zuzana Roth, ``Inequalities for entropy,
  Hausdorff dimension, and Lipschitz constants,''} \emph{Studia Math.}
  \textbf{250} (2020), 253--264; arXiv:1807.01196. Defines the HausLip
  constant \[
  \mathrm{HausLip}(X,f)=\inf_{d\in\mathcal D(X)}\mathrm{HD}_d(X)\log^+\operatorname{Lip}_d(f),
  \] proves it is a topological conjugacy invariant and \(\ge h(f)\),
  and proves \(\mathrm{HausLip}(X,f)=h(f)\) for positively expansive
  maps and for linear maps of tori. The proof for expansive maps
  constructs a metric in which
  \(\operatorname{Lip}_D(f^n)\le16\,\mathrm{SkewLoc}_D(f^n)\), i.e.~a
  metric that is uniformly expanding up to a bounded factor. This is the
  closest published construction to your visual-metric theorem.
\item
  \textbf{S. A. Dovbysh, ``Optimal Lyapunov metrics of expansive
  homeomorphisms,''} \emph{Izv. Math.} \textbf{70} (2006), 883--929; DOI
  10.1070/IM2006v070n05ABEH002332. Sharpens Reddy, Sakai, and Fried:
  every expansive homeomorphism of a metrizable compactum admits a
  compatible Lyapunov metric; for A$^{\#}$-homeomorphisms the lower
  bounds for the contraction and expansion constants are attained
  simultaneously by one metric satisfying Ruelle's Lipschitz axiom. This
  is the two-sided analogue of optimizing the metric in your
  \(\lambda_{\exp}\).
\end{enumerate}

\begin{center}\rule{0.5\linewidth}{0.5pt}\end{center}

\section{Your objects in the literature's
language}\label{your-objects-in-the-literatures-language}

{\def\LTcaptype{none} % do not increment counter
\begin{longtable}[]{@{}
  >{\raggedright\arraybackslash}p{(\linewidth - 4\tabcolsep) * \real{0.3333}}
  >{\raggedright\arraybackslash}p{(\linewidth - 4\tabcolsep) * \real{0.3333}}
  >{\raggedright\arraybackslash}p{(\linewidth - 4\tabcolsep) * \real{0.3333}}@{}}
\toprule\noalign{}
\begin{minipage}[b]{\linewidth}\raggedright
Draft object
\end{minipage} & \begin{minipage}[b]{\linewidth}\raggedright
Closest literature object
\end{minipage} & \begin{minipage}[b]{\linewidth}\raggedright
Main references
\end{minipage} \\
\midrule\noalign{}
\endhead
\bottomrule\noalign{}
\endlastfoot
\((\lambda,\epsilon)\)-expanding & expanding in the sense of Ruelle;
distance-expanding; positively expansive + expanding metric & Reddy
1982; Przytycki--Urba\'{n}ski 2010; Shub 1969 \\
\(\lambda_{\exp}(f)\) & maximal expansion constant; logarithmic
expansion rate \(E\); expanding ratio & Hirsch 1994; Schreiber 1997;
Fujita--Kato--Matsumoto 2010; Dovbysh 2006 \\
\(d(x,y;\mathcal U)\) & chain distance in a cover; uniformity chain &
Bowen 1971/1973; Misiurewicz 1986; uniform-space entropy \\
\(h_{\mathrm{dist}}(x,y)\) & pairwise Lyapunov exponent / pointwise
expansion rate & Oseledets; Hirsch 1994; Schreiber 1997 \\
\(r(\mathcal U;\mathcal V)\) & reciprocal of the Lebesgue number & Sun
2012 \\
\(h_{\min}\) & minus the Lebesgue-number decay rate; infimum Lyapunov
exponent (smooth case) & Sun 2012; Schreiber 1997 \\
\(h_{\mathrm{diam}}\), \(h_{\max}\) & diameter growth; maximal
characteristic exponent;
\(R(f)=\lim\frac1n\log\sup_x\lVert Df^n_x\rVert\) & Kifer 1983; Newhouse
1988; Yomdin 1987; Pacifico--Vieitez 2020 \\
\(\lambda_{\mathrm{Lip}}\) & infimum of Lipschitz constants over length
metrics; conjugacy-class Lipschitz infimum; HausLip constant &
Bobok--Roth 2019; Roth--Roth 2020; Milnor Section 13 \\
\(h_{\mathrm{top}}\) & topological entropy; volume growth &
Adler--Konheim--McAndrew 1965; Bowen 1971; Newhouse 1988; Yomdin 1987 \\
\end{longtable}
}

One more caution on the table: the draft's \(\lambda_{\mathrm{Lip}}\) is
an infimum over \emph{length} metrics, while the Bobok--Roth and
Roth--Roth literature usually optimizes over all compatible metrics or
over Euclidean metrics in a conjugacy class. The quantities are closely
related but not literally the same; the table is a map of ideas, not an
identity of definitions. Section 8 takes the maximum-end rows up in
detail, including why the length-metric restriction is not optional.

Three terminology warnings.

\begin{itemize}
\item
  \textbf{``Expansive'' vs ``positively expansive.''} In the literature,
  \emph{expansive} usually means the two-sided homeomorphism notion:
  \(d(f^n x,f^n y)<\gamma\) for all \(n\in\mathbb Z\) implies \(x=y\).
  Your \(\mathbb N\)-expansive is \emph{positively expansive}. For a
  homeomorphism of a compact space, positive expansiveness is extremely
  restrictive (it forces a finite space; see the note ``Positively
  expansive homeomorphisms of compact spaces,'' arXiv:math/0308294, and
  the references there). Your one-sided setting is the right one for
  non-invertible maps, but the word ``expansive'' alone will be read as
  the two-sided notion by many dynamicists; the draft already writes
  ``\(\mathbb N\)-expansive or just expansive,'' which is a good guard.
\item
  \textbf{``Expansive constant'' vs ``expansion constant.''} An
  \emph{expansive constant} is usually the \textbf{scale} \(\epsilon\)
  (Sun 2011, Walters, Gottschalk--Hedlund). An \emph{expansion constant}
  / \emph{expanding ratio} / \emph{expansion rate} is the \textbf{rate}
  \(\lambda\) (Hirsch, Schreiber, Fujita--Kato--Matsumoto). Your
  \(\lambda_{\exp}\) is a rate; your \(\epsilon\) is a scale. The phrase
  ``maximal expanding constant'' is ambiguous between the two, and the
  literature uses both.
\item
  \textbf{The refinement order runs the other way in your draft.} You
  write \(\mathcal U\le\mathcal V\) when \(\mathcal U\) refines
  \(\mathcal V\), so your \(\vee\) is the common coarsening and your
  \(\wedge\) the common refinement. In dynamics the opposite order is
  universal: refinement makes a cover \emph{larger}, the common
  refinement is the join, and \(N\), \(H\) and \(h(\cdot;\mathcal U)\)
  all increase under refinement. Adler--Konheim--McAndrew write
  \(\mathfrak{A}\prec\mathfrak{B}\) for ``\(\mathfrak{B}\) refines
  \(\mathfrak{A}\)'', and their Property 3 fixes the direction without
  any guesswork: they prove
  \(\mathfrak{A}<\mathfrak{B}\Rightarrow H(\mathfrak{A}\vee\mathfrak{B})=H(\mathfrak{B})\),
  and the common refinement has the entropy of the \emph{finer} cover,
  so \(\mathfrak{B}\) is the finer one and \(<\) points from coarse to
  fine. Sun (2012) writes \(\mathcal U\succ\mathcal V\) for
  ``\(\mathcal U\) refines \(\mathcal V\)'' (his Lemma 2.2: refinement
  decreases the Lebesgue number); Bruin's notes write
  \(\mathcal U\succeq\mathcal V\); Walkden's notes write
  \(\alpha\le\beta\) for ``\(\beta\) refines \(\alpha\)''; the
  mean-dimension literature writes \(\beta\succ\alpha\) for ``\(\beta\)
  refines \(\alpha\)''. The order with refinement \emph{decreasing} a
  cover, and the common refinement as the meet, is the one used in
  covering uniformities (Willard; Isbell: a uniformity is closed under
  coarsening, and its meet is the common refinement) --- not the one
  used in dynamics. The notion of refinement itself is of course the
  same everywhere (Engelking, Section 1.5; Willard, Section 36); only its order
  differs. Two practical consequences. First, the glyph does not
  identify the order: \(\prec\) means ``refines'' in general topology
  and ``is coarser than'' in Adler--Konheim--McAndrew, and \(\le\)
  appears in both directions in the sources above, so a paper has to say
  which order it uses. Second, once the order is declared the rest
  follows from the standard lattice meanings --- \(\wedge\) is the meet,
  \(\vee\) the join --- so nothing further needs saying, and your draft
  does declare it: \texttt{covers.tex} writes
  \(\mathcal V\le\mathcal U\) for ``\(\mathcal V\) refines
  \(\mathcal U\)'', and its remark after the definition of
  \(\mathcal U\wedge\mathcal V\) records the reversal from AKM together
  with the corresponding reversal of meet and join. The reader's whole
  task is then to apply the lattice meanings to that order: the common
  refinement, which dynamics writes \(\vee\), is the \emph{meet} here
  and appears as \(\wedge\) --- it is the operation in
  \(\mathcal U_n=\mathcal U\wedge f^{-1}\mathcal U\wedge\cdots\wedge f^{-(n-1)}\mathcal U\),
  and \(\vee\) itself never occurs in the paper outside that remark. The
  convention is also consistent at the empty end: \(\mathcal U_0=\{X\}\)
  is the coarsest cover, which is what an empty meet should be in this
  order.
\end{itemize}

\textbf{Where the two orders come from.} Adler--Konheim--McAndrew's
footnote to ``reflexive partial ordering'' reads ``For terminology see
J. L. Kelley, \emph{General topology}, Van Nostrand, New York, 1955'',
so they present the terminology as inherited. But Kelley attaches no
inequality to refinement of \emph{covers} --- he states the relation in
words --- and his own covering-uniformity condition (Problem J, p.~207)
is closed under coarsening, which is the orientation topology kept. The
only refinement order Kelley does symbolize is for subdivisions in the
chapter on integration, where \(S\ge S'\) iff \(S\) refines \(S'\): the
direction AKM use. What AKM's direction matches is the \emph{partition}
convention of ergodic theory, and their paper says so: the common
refinement of two covers is called the join (Definition 2), the remark
after Property 1 has \(\mathfrak{A}<\mathfrak{A}\vee\mathfrak{B}\) and
\(\mathfrak{B}<\mathfrak{A}\vee\mathfrak{B}\) and concludes that the
covers form a directed set, and in Section 4 the same \(\vee\) is used for
measurable partitions, ``the common refinement'', with
\(H_\mu(\mathfrak{A}\vee\mathfrak{B})\le H_\mu(\mathfrak{A})+H_\mu(\mathfrak{B})\).
So the dynamics order is AKM's own alignment of covers with partitions;
it is not a misreading of Kelley, and it is not the covering-uniformity
order. The partition convention itself is older than the covers paper
--- it comes from the founding papers of metric entropy, Kolmogorov's
1958 note and Sinai's 1959 note --- and is written out explicitly in
Rokhlin's survey (Section 1.3): he fixes \(\xi<\xi'\) for ``\(\xi'\) is a
subpartition of \(\xi\)'', defines the product
\(\textstyle\bigvee_\alpha\xi_\alpha\) as the join (every \(\xi_\alpha\)
is below it, and it is the least such partition) and the intersection
\(\textstyle\bigwedge_\alpha\xi_\alpha\) as the meet, with the point
partition \(\varepsilon\) at the top and the trivial partition \(\nu\)
at the bottom. So the measure-theoretic convention and AKM's cover
convention were already the same one when AKM wrote: their paper was
received in August 1963, and its first definition says it is conforming
to ergodic theory. (The two Doklady notes are not on mathnet and are
behind a paywall, so I have not read them; the earliest statement in
print I have seen is Rokhlin's.)

\subsection{Three kinds of invariant, and which is
which}\label{three-kinds-of-invariant-and-which-is-which}

Three kinds of object appear in this report, and they should not be
conflated:

{\def\LTcaptype{none} % do not increment counter
\begin{longtable}[]{@{}
  >{\raggedright\arraybackslash}p{(\linewidth - 4\tabcolsep) * \real{0.3333}}
  >{\raggedright\arraybackslash}p{(\linewidth - 4\tabcolsep) * \real{0.3333}}
  >{\raggedright\arraybackslash}p{(\linewidth - 4\tabcolsep) * \real{0.3333}}@{}}
\toprule\noalign{}
\begin{minipage}[b]{\linewidth}\raggedright
Kind
\end{minipage} & \begin{minipage}[b]{\linewidth}\raggedright
Invariance
\end{minipage} & \begin{minipage}[b]{\linewidth}\raggedright
Examples in this report
\end{minipage} \\
\midrule\noalign{}
\endhead
\bottomrule\noalign{}
\endlastfoot
\textbf{Topological} & invariant under topological conjugacy &
\(h_{\mathrm{top}}\); your \(h_{\min}\), \(h_{\mathrm{dist}}\),
\(h_{\mathrm{diam}}\), \(h_{\max}\); your \(\lambda_{\exp}\) and
\(\lambda_{\mathrm{Lip}}\); \(\mathrm{HausLip}\) (Roth--Roth);
\(\Lambda(f)\) (Bobok--Roth); \(\lambda(f)\) (Braga--Sieg) \\
\textbf{Metric-bound} & invariant under bi-Lipschitz changes of the
metric, not under topological conjugacy & Sun's \(h_L^\pm\) and
\(h_E^\pm\); Kifer's \(\Lambda^\pm(x)\), \(\lambda^\pm(x)\) in a fixed
metric; Morales--Thieullen--Villavicencio's \(\chi_f\); Hirsch's
one-step \(EC(f,K)\) \\
\textbf{Smooth} & needs a differentiable structure; independent of the
choice of Riemannian metric & the Lyapunov exponents \(\lambda_i\);
Schreiber's \(E(f,K)\); Newhouse--Yomdin's \(v_k(f)\) and \(R(f)\); the
Lyapunov spectrum \([\alpha^-,\alpha^+]\); Bowen's equation and the
pressure formula \\
\end{longtable}
}

\textbf{Every invariant in your draft is topological}, and the reason is
the same in each case: it is built either from open covers (the \(h\)'s
and \(h_{\mathrm{dist}}\)) or as a supremum or infimum over \emph{all}
compatible metrics (\(\lambda_{\exp}\), \(\lambda_{\mathrm{Lip}}\)). A
homeomorphism carries the class of compatible metrics to itself, so a
supremum over that class is automatically invariant under topological
conjugacy. That is a design decision, not an accident, and it is what
separates \(\lambda_{\exp}\) from every smooth expansion rate in Section
4. The distinction is precisely \emph{all compatible metrics} versus
\emph{all Riemannian metrics}: the first supremum is topological, the
second is a smooth invariant --- Schreiber's \(E\) (Section 4.4). The
two are suprema over different classes, but the inclusion does give one
inequality, as you pointed out: a Riemannian metric in which \(f^n\)
expands by \(\rho\) yields a \emph{compatible} metric in which \(f\)
expands by \(\rho^{1/n}\), so \(\lambda_{\exp}\ge e^{E}\) (the one-line
proof is in Section 4.4). The reverse inequality is the open half; by
your theorem it is the same as \(h_{\min}\le E\).

\textbf{Most of the literature's quantities are not topological.} Sun's
\(h_L\) is invariant under bi-Lipschitz changes of the metric and not
under topological conjugacy (his Examples 6.4--6.5); his \(h_E\) is
metric-bound too, since it is read off a separation scale in a fixed
metric; Kifer's exponents and Morales--Thieullen--Villavicencio's
\(\chi_f\) are bi-Lipschitz invariant only (Section 8.4); the Lyapunov
exponents, Schreiber's \(E\), and Newhouse--Yomdin's \(v_k\) and \(R\)
need a differentiable structure. So ``your \(h_{\min}\) is Sun's
\(h_L\)'' is a comparison between a topological invariant and a
metric-bound one, and the honest statement is the comparison
\(h_{\min}\approx-h_L\) \emph{in a fixed metric and for an expansive
cover}, not an identity of invariants.

\textbf{The interesting theorems are the bridges between the columns.}

\begin{itemize}
\tightlist
\item
  \emph{smooth to smooth}: Schreiber's \(E(f,K)=\inf_\mu E(f,\mu)\)
  identifies Hirsch's asymptotic metric expansion rate with the infimum
  of the Lyapunov exponents.
\item
  \emph{topological to metric-bound}: Reddy's positively expansive
  \(\Leftrightarrow\) expands small distances in some metric, and
  Fathi's hyperbolic metric for expansive homeomorphisms.
\item
  \emph{topological to metric}: the entropy--dimension--expansion
  inequalities (Fathi, Dai--Zhou--Geng, Misiurewicz,
  Fujita--Kato--Matsumoto, Roth--Roth), which relate the topological
  \(h_{\mathrm{top}}\) to a dimension in a chosen metric times a
  logarithmic expansion or Lipschitz rate.
\item
  \emph{topological to topological}: your theorem
  \(\log\lambda_{\exp}=h_{\min}\) --- both sides are topological, and
  the content is that a supremum over metrics equals a cover-chain
  growth rate.
\end{itemize}

Each subsection below says which kind it is discussing.

\begin{center}\rule{0.5\linewidth}{0.5pt}\end{center}

\section{Theme 1: entropy-like invariants built from
distances}\label{theme-1-entropy-like-invariants-built-from-distances}

\subsection{Entropy from separated and spanning sets (Bowen,
Dinaburg)}\label{entropy-from-separated-and-spanning-sets-bowen-dinaburg}

The standard distance-based definition of topological entropy is \[
d_n^f(x,y)=\max_{0\le k<n}d(f^k x,f^k y),\qquad
h_{\mathrm{top}}(f)=\lim_{\epsilon\to0}\limsup_{n\to\infty}\frac1n\log N(n,\epsilon),
\] where \(N(n,\epsilon)\) is the maximal cardinality of an
\((n,\epsilon)\)-separated set (or the minimal cardinality of an
\((n,\epsilon)\)-spanning set). This is ``distance'' in the sense that
it uses the metric, but it is still a \textbf{volume} count: it counts
how many \(\epsilon\)-separated orbit segments there are. The metric
enters the definition, but the \(\epsilon\to0\) limit removes it:
\(h_{\mathrm{top}}\) is topological, and this is the one place in Theme
1 where a metric-based definition and a topological outcome agree
exactly.

Your \(d(x,y;\mathcal U^{(n)})\) is the cover analogue of
\(d_n^f(x,y)\), and your \(h_{\mathrm{dist}}(x,y;\mathcal U)\) is the
exponential growth rate of a \textbf{single pair's} \(d_n^f\)-distance,
not the count of separated sets. The standard entropy is the sup over
scales of the growth rate of the number of pairs; your
\(h_{\mathrm{dist}}\) is the growth rate of one pair. The two are
related by the variational principle and by the Lyapunov spectrum (see
Section 3.7), but they are not the same construction.

\subsection{\texorpdfstring{Entropy from \(\epsilon\)-chains
(Misiurewicz
1986)}{Entropy from \textbackslash epsilon-chains (Misiurewicz 1986)}}\label{entropy-from-epsilon-chains-misiurewicz-1986}

A genuinely chain-based definition of entropy is due to Misiurewicz. For
\(\delta,\epsilon>0\), let \(S(n,\delta,\epsilon)\) be the maximal
cardinality of a set of \(\delta\)-separated \(\epsilon\)-chains of
length \(n\) (an \(\epsilon\)-chain is a sequence \(x_0,\dots,x_{n-1}\)
with \(d(fx_i,x_{i+1})\le\epsilon\); two chains are \(\delta\)-separated
if \(d(x_i,y_i)\ge\delta\) for some \(i\)). Then \[
h(f)=\lim_{\delta\to0}\lim_{\epsilon\to0}\limsup_{n\to\infty}\frac1n\log S(n,\delta,\epsilon).
\] This is the closest classical ``chain'' definition to your
cover-chain construction, but it counts chains rather than tracking the
chain distance of one pair. It shows that chain distances can recover
the volume entropy, which is a useful sanity check for any chain-based
definition. A modern account with this formula is Richeson--Wiseman
(2008). The formula is topological as well: the \(\delta\to0\),
\(\epsilon\to0\) limits remove the metric, so this is a chain formula
for the topological invariant, not a metric-bound quantity.

\subsection{\texorpdfstring{Entropy from Lebesgue numbers (Sun 2012):
the closest invariant to
\(h_{\min}\)}{Entropy from Lebesgue numbers (Sun 2012): the closest invariant to h\_\{\textbackslash min\}}}\label{entropy-from-lebesgue-numbers-sun-2012-the-closest-invariant-to-h_min}

For an open cover \(\mathcal U\) of a compact metric space, let
\(\delta(\mathcal U)\) be its Lebesgue number: the largest \(\delta>0\)
such that every ball of radius \(\le\delta\) is contained in some
element of \(\mathcal U\). Put
\(\delta_n(f,\mathcal U)=\delta(\bigvee_{k<n}f^{-k}\mathcal U)\), and
define \[
h_L^-(f,\mathcal U)=\liminf_{n\to\infty}-\frac1n\log\delta_n(f,\mathcal U),\qquad
h_L^+(f,\mathcal U)=\limsup_{n\to\infty}-\frac1n\log\delta_n(f,\mathcal U),
\] and \(h_L^\pm(f)=\sup_{\mathcal U}h_L^\pm(f,\mathcal U)\) over finite
open covers.

Sun proves:

\begin{itemize}
\tightlist
\item
  \(h_L^\pm(f)\ge0\), \(h_L^-(f)\le h_L^+(f)\);
\item
  \(h_L^\pm(f^n)=n\,h_L^\pm(f)\);
\item
  if \(\mathcal U\) refines \(\mathcal V\) (in Sun's notation
  \(\mathcal U\succ\mathcal V\)), then
  \(h_L^\pm(f,\mathcal U)\ge h_L^\pm(f,\mathcal V)\);
\item
  \(\dim_B(X)\,h_L^-(f)\ge h(f)\) and \(\dim_H(X)\,h_L^+(f)\ge h(f)\);
\item
  if \(f\) is Lipschitz with constant \(L(f)\), then
  \(h_L^+(f,\mathcal U)\le\log L(f)\) for every \(\mathcal U\);
\item
  \(h_L^\pm\) is invariant under bi-Lipschitz conjugacy, but
  \textbf{not} a topological invariant (Examples 6.4--6.5 give
  topologically conjugate homeomorphisms with different \(h_L\));
\item
  if \(f\) is an expansive homeomorphism with expansive constant
  \(\gamma\), then every open cover of diameter \(<\gamma\) computes
  \(h_L^\pm(f)\) (Corollary 3.9).
\end{itemize}

\textbf{Comparison with your in-radius.} Your
\(r(\mathcal U;\mathcal V)\) is the largest \(k\) such that every
\(\mathcal V\)-chain of length \(\le k\) lies in a single
\(\mathcal U\)-element. If \(\mathcal V\) has mesh \(\eta\) and
\(\mathcal U\) has Lebesgue number \(\delta\), then a
\(\mathcal V\)-chain of length \(k\) has diameter at most \(k\eta\), so
\(k\eta<\delta\) implies \(r(\mathcal U;\mathcal V)\ge k\); conversely,
a maximal failing chain gives a lower bound of the same shape. Thus \[
r(\mathcal U;\mathcal V)\asymp \frac{\delta(\mathcal U)}{\operatorname{mesh}(\mathcal V)}
\] up to constants, and for the dynamical sequence
\(\mathcal V=\mathcal U^{(n)}\), \[
h_{\min}=\lim_{n\to\infty}\frac1n\log r_{\mathcal U}(n)
\;\approx\;-\lim_{n\to\infty}\frac1n\log\delta_n(f,\mathcal U)
=-h_L(f,\mathcal U).
\] The approximation is not an identity: it uses that, for an expansive
cover in a compatible metric, the mesh and the Lebesgue number of
\(\mathcal U^{(n)}\) decay at the same exponential rate. That is exactly
the situation in which your theorem makes \(h_{\min}\) independent of
\(\mathcal U\). Sun's Corollary 3.9 is the analogous statement for
\(h_L\) and expansive homeomorphisms.

This is, as far as I found, the closest prior work to your \(h_{\min}\).
The main difference is that Sun works with a fixed metric and a
one-cover Lebesgue number, and his \(h_L\) is \textbf{not} a topological
invariant but a metric-bound one, invariant under bi-Lipschitz changes
of the metric only; your \(h_{\min}\) is defined by minimizing over
pairs and suping over covers, and your theorem says it \emph{is} a
topological invariant for expansive maps. In the language of the draft,
Sun's \(h_L\) is the metric-bound upper/lower version of
\(-\log\lambda_{\exp}\); passing from \(h_L\) to \(h_{\min}\) is exactly
the passage from a metric invariant to a topological one.

\subsection{Entropy from expansive constants (Sun
2011)}\label{entropy-from-expansive-constants-sun-2011}

A related but distinct ``distance'' invariant is the exponential decay
rate of the \textbf{expansive constant}. For an expansive map \(f\), let
\(\gamma(f)\) be the largest \(\epsilon>0\) such that
\(d(f^n x,f^n y)<\epsilon\) for all \(n\ge0\) implies \(x=y\). Define \[
h_E^+(f)=\limsup_{n\to\infty}-\frac1n\log\gamma(f^n),\qquad
h_E^-(f)=\liminf_{n\to\infty}-\frac1n\log\gamma(f^n).
\] Sun proves \(h_E^\pm(f)\ge0\), \(h_E^-(f)\le h_E^+(f)\),
\(h_E^\pm(f)=h_E^\pm(f^{-1})\) for homeomorphisms,
\(h_E^+(f^n)\le n h_E^+(f)\), \(h_E^-(f^n)\ge n h_E^-(f)\),
\(h_E^\pm(f)\le\log L(f)\) for Lipschitz \(f\), and the dimension bound
\(h_E^-(f)\,\dim_B^+X\ge h(f)\), \(h_E^+(f)\,\dim_B^-X\ge h(f)\). For a
two-sided subshift of finite type with metric
\(d(\omega,\omega')=q^{-\min\{|j|:\omega_j\ne\omega'_j\}}\),
\(h_E^\pm(\sigma_A)=\frac12\log q\), so
\(h_E\cdot\dim_H=h_{\mathrm{top}}\); the one-sided case is analogous
with the corresponding dimension.

The expansive constant is a \textbf{separation scale}, not an in-radius:
it is the largest scale at which distinct orbits are eventually
separated, whereas your \(r_{\mathcal U}(n)\) is the largest number of
steps that fit inside a cover element. The two are related (both measure
how fast the dynamics resolves the space), but they are not the same
invariant. Sun's inequality \(h_E\le h_L\) is a relation between the two
distance-based quantities. Both are metric-bound, not topological: each
is read off a scale in a fixed metric and changes only by a bounded
amount under a bi-Lipschitz change of that metric.

\subsection{Distance entropy (Dai--Jiang 2008): a different notion with
the same
name}\label{distance-entropy-daijiang-2008-a-different-notion-with-the-same-name}

The literature already uses the phrase \textbf{distance entropy} for a
different construction. For a separable metric space \(X\), a continuous
\(f:X\to X\), a compact \(K\subseteq X\) and \(\epsilon>0\), the
\emph{step length} \(n(f,K,\epsilon)\) is the largest \(n\) such that
\(\operatorname{diam}(f^j(K))<\epsilon\) for all \(0\le j<n\); set
\(\operatorname{diam}_\epsilon(K)=e^{-n(f,K,\epsilon)}\). Then a
Hausdorff-type measure is built from covers by sets with prescribed step
length, and its critical exponent is \(\operatorname{ent}_H(f;Y)\), the
distance entropy of \(f\) on \(Y\). Dai and Jiang prove that if \(f\) is
locally Lipschitz with constant \(l\), then
\(\operatorname{ent}_H(f;Y)\le\max\{0,\mathrm{HD}(Y)\log l\}\), and if
\(f\) is locally expanding with skewness \(\lambda\), then
\(\operatorname{ent}_H(f;Y)\ge\mathrm{HD}(Y)\log\lambda\); they also
prove countable stability.

This is a \textbf{Hausdorff-dimension} construction using the time a set
stays small, not the growth rate of a pair's distance. It is a valid
``other choice'' of how to turn distances into an entropy, but it is not
your \(h_{\mathrm{dist}}\). The name collision is worth flagging in the
draft if you use ``distance entropy'' for \(h_{\mathrm{dist}}\).
Dai--Jiang's invariant is metric-bound in the same way --- Hausdorff
dimension and step lengths in a fixed metric --- so it is not a
topological invariant of \(f\).

\subsection{Length and volume growth (Newhouse,
Yomdin)}\label{length-and-volume-growth-newhouse-yomdin}

The phrase ``length growth rather than volume growth'' in your abstract
has a precise smooth analogue. For a \(C^1\) map \(f\) of a compact
Riemannian manifold and \(0\le k\le\dim M\), let \[
v_k(f)=\sup_{D}\limsup_{n\to\infty}\frac1n\log\operatorname{vol}_k(f^n(D)),
\] the supremum over \(k\)-dimensional disks \(D\). Newhouse proved
\(v_k(f)\ge h_{\mathrm{top}}(f)\) for \(C^r\) maps with \(r>1\); Yomdin
proved \(v_k(f)\le h_{\mathrm{top}}(f)+k\,R(f)\) for \(C^r\) maps, where
\(R(f)=\lim_n\frac1n\log\sup_x\|df^n_x\|\), and hence
\(v_\infty(f)=h_{\mathrm{top}}(f)\) for \(C^\infty\) maps. For \(k=1\)
this is the exponential growth rate of lengths of curves; for
\(k=\dim M\) it is the volume growth.

The point for your draft is that \textbf{distance growth is the \(k=0\)
end of a dimension-indexed family}, with entropy at the \(k=\dim M\)
end. Your \(h_{\mathrm{dist}}(x,y)\) is a pointwise \(k=0\) growth rate;
your \(h_{\mathrm{diam}}\) and \(h_{\max}\) are still \(k=0\)-type
quantities (diameters and pair distances); \(h_{\mathrm{top}}\) is the
\(k=\dim M\) volume growth. The Newhouse--Yomdin family gives a natural
interpolation, and the ``distance vs volume'' dichotomy is the two ends
of it. If you want an ``other choice'' that is intermediate between your
\(h_{\max}\) and \(h_{\mathrm{top}}\), the \(k\)-volume growth rates are
the natural candidates. The whole family is smooth: \(v_k\) and \(R\)
are built from the derivative and the Riemannian volume, they are
independent of the choice of metric, and they are defined only for
smooth maps of manifolds (Newhouse's inequality needs
\(C^{1+\epsilon}\), Yomdin's \(C^k\)). The comparison with your
\(h_{\mathrm{dist}}\), \(h_{\mathrm{diam}}\), \(h_{\max}\) is an analogy
between a smooth family and a topological one, not an identity.

A related ``volume'' notion is the \textbf{expansion entropy} of Hunt
and Ott: the exponential growth rate of the average of the maximal
\(d\)-dimensional volume growth ratio along orbits, minus the rate at
which volume leaves the restraining region. On an invariant set it is
the sum of the positive Lyapunov exponents; for \(C^\infty\) maps it
equals topological entropy (by Kozlovski). It generalizes
Sacksteder--Shub's entropy and is the volume/Lyapunov-sum analogue of
your distance-based expansion.

\subsection{Lyapunov exponents and the growth
spectrum}\label{lyapunov-exponents-and-the-growth-spectrum}

Your \(h_{\mathrm{dist}}(x,y)\) is the exponential growth rate of the
distance between the orbits of \(x\) and \(y\). For a \(C^1\) map, this
is a \textbf{pairwise Lyapunov exponent}: the growth rate of
\(d(f^n x,f^n y)\) is governed by the largest Lyapunov exponent of the
pair (generically the larger of the two exponents, unless the pair is
trapped in a stable direction). The minimum over pairs of this rate is
the smallest Lyapunov exponent; the maximum is the largest. Thus your \[
h_{\min}\le h_{\mathrm{diam}}\le h_{\max}\le h_{\mathrm{top}}
\] is the topological analogue of \[
\lambda_{\min}\le \text{(diameter growth)}\le \lambda_{\max}\le \sum_{\lambda_i>0}\lambda_i
\] for smooth systems, with the last term being the volume growth
(Pesin/Ruelle). The middle term depends on which notion of ``diameter''
you use; for a linear map it is the largest singular value, for a
conformal map it is the common expansion rate.

The \textbf{Lyapunov spectrum} literature studies the level sets of the
pointwise exponents: for \(\alpha\), the set of points with exponent
\(\alpha\), its Hausdorff dimension, and the multifractal spectrum. Your
``growth spectrum'' is the set of values of \(h_{\mathrm{dist}}(x,y)\);
the natural refinement is the function \(\alpha\mapsto\) (size of the
set of pairs with \(h_{\mathrm{dist}}=\alpha\)). I did not find a paper
that studies this for the cover-chain distance in the topological
setting, but the smooth analogues are the multifractal analysis of
Lyapunov exponents (Barreira--Pesin--Schmeling; Pesin--Weiss;
Iommi--Jordan for interval maps). This is a possible ``other choice'' if
you want the spectrum to be a function rather than a set. This whole
paragraph is smooth: the exponents need a derivative and are not
topological invariants. What survives topologically is the sign ---
positivity of the exponents, and the existence of a metric in which all
of them are nonzero for an expansive homeomorphism (Fathi's metric;
Section 8.4).

\subsection{Visual metrics and conformal
dimension}\label{visual-metrics-and-conformal-dimension}

Your visual graph construction is a Floyd-style version of the
\textbf{visual metric} on the boundary of a Gromov hyperbolic space. The
Gromov product \(\langle x,y\rangle_{\mathcal U}\) is the largest \(n\)
such that \(x,y\) are \(\mathcal U^{(n)}\)-connected, and the visual
metric is \(\lambda^{-\langle x,y\rangle}\) for a discount factor
\(\lambda\). The standard references are M. Gromov, ``Hyperbolic
groups,'' in \emph{Essays in Group Theory}, MSRI Publ. 8, Springer 1987,
75--263; M. Bonk and O. Schramm, ``Embeddings of Gromov hyperbolic
spaces,'' \emph{Geom. Funct. Anal.} \textbf{10} (2000), 266--306; S.
Buyalo and V. Schroeder, \emph{Elements of Asymptotic Geometry}, EMS
2007. Floyd's original construction is W. J. Floyd, ``Group completions
and limit sets of Kleinian groups,'' \emph{Invent. Math.} \textbf{57}
(1980), 205--218. A note on kind: everything in this subsection is
metric or conformal. Visual metrics depend on the discount factor, and
\(\mathrm{ARCdim}\) is invariant under quasisymmetric changes of metric
but not under topological conjugacy.

The closest dynamical use is Ha\"{i}ssinsky--Pilgrim's \textbf{coarse
expanding conformal} (cxc) theory. For a topologically cxc map of degree
\(d\), they construct visual metrics \(d_\epsilon\) in which small balls
are sent homeomorphically and homothetically onto their images with
expansion constant \(e^\epsilon\), and the Hausdorff dimension is
\(q=\log d/\epsilon\). The \textbf{Ahlfors regular conformal dimension}
\(\mathrm{ARCdim}\) is the infimum of Hausdorff dimensions over the
conformal gauge; in the cxc setting it corresponds to the supremum of
the admissible expansion exponents \(\epsilon\), i.e.~to
\(\log d/\mathrm{ARCdim}\). Your \(\lambda_{\exp}\) is a supremum over
\textbf{all} compatible metrics, not just the conformal gauge, so
\(\lambda_{\exp}\ge e^{\epsilon_{\max}}\) in the cxc case; whether
equality holds is a separate question (and not one to settle here).

The thermodynamic formalism gives the general version of the equality:
for a conformal expanding repeller, the Hausdorff dimension \(s\) of the
limit set is the root of Bowen's equation \(P(-s\log|f'|)=0\), where
\(P\) is the topological pressure (Bowen 1979; Ruelle 1982;
Parry--Pollicott 1990). This is the standard ``distance vs volume''
machine in the conformal case: the derivative (distance expansion)
enters as a potential, and the pressure (volume growth) determines the
dimension.

The energy side is your own paper with Pilgrim: K. M. Pilgrim and D. P.
Thurston, ``Ahlfors-regular conformal dimension and energies of graph
maps,'' arXiv:2112.09041, \emph{Groups Geom. Dyn.} (2024/2025), DOI
10.4171/ggd/914. The asymptotic \(q\)-conformal energy
\(\overline E^q[\pi,\phi]\) is continuous and non-increasing in \(q\);
\(\mathrm{ARCdim}\) lies in the interval where \(\overline E^q=1\). In
the draft's comb-models section, the asserted identity
\(h_{\min}=-\log\overline E^\infty[\pi,\phi]\) is the \(q=\infty\) end
of this family, i.e.~the maximal expansion rate, while
\(\mathrm{ARCdim}\) is the \(q\) where the energy crosses 1. The two are
different invariants, and it is worth keeping them explicitly separate:
the maximal expanding constant is the reciprocal of the
\(\infty\)-energy / stretch factor, not the reciprocal of the conformal
dimension. (This is the user's own work; I am flagging the distinction,
not claiming a new result. The draft's bibliography key for
Ha\"{i}ssinsky--Pilgrim is \texttt{HP12:AlgExpanding}.)

\begin{center}\rule{0.5\linewidth}{0.5pt}\end{center}

\section{Theme 2: the maximal expanding
constant}\label{theme-2-the-maximal-expanding-constant}

\subsection{Existence and equivalent
definitions}\label{existence-and-equivalent-definitions}

The standard notion is: \(f\) is \textbf{expanding} if there are
\(\lambda>1\) and \(\epsilon>0\) such that \(d(fx,fy)\ge\lambda d(x,y)\)
whenever \(d(x,y)<\epsilon\). Equivalent formulations in the literature:

\begin{itemize}
\tightlist
\item
  \textbf{Positively expansive + expanding metric.} Reddy proved that a
  map of a compact metric space is positively expansive iff it expands
  small distances iff it increases small distances. So every positively
  expansive map has an expanding metric.
\item
  \textbf{Ruelle expanding.} A continuous open surjection that expands
  distances locally. On compact manifolds, positive expansiveness
  implies openness (invariance of domain), so the notions agree there.
\item
  \textbf{Distance-expanding.} There are \(\lambda>1\) and \(C>0\) such
  that \(\rho(f^n x,f^n y)\ge C\lambda^n\rho(x,y)\) for all \(n\) and
  all \(x,y\) sufficiently close (Przytycki--Urba\'{n}ski,
  \emph{Conformal Fractals}, Ch. 4). This is the form used in
  thermodynamic formalism. Of the three, \textbf{positive expansiveness
  is topological} --- it does not depend on the compatible metric ---
  while \textbf{distance-expanding is metric-bound}: it is a property of
  one fixed metric. Ruelle's formulation sits between them: the openness
  is topological, the local expansion is metric. ``Expanding in some
  compatible metric'' is therefore a topological property even though
  ``expanding in this metric'' is not (Reddy, Section 2.1).
\end{itemize}

\subsection{Fixed metric: the expansion constant and Hirsch's
rate}\label{fixed-metric-the-expansion-constant-and-hirschs-rate}

For a map \(f\) of a compact metric space and a compact positively
invariant set \(K\), the \textbf{expansion constant} is \[
EC(f,K)=\sup\{\rho\ge0:\ \forall \rho^*<\rho\ \exists\delta>0,\ d(x,y)\le\delta,\ x\in K\implies d(fx,fy)\ge\rho^*d(x,y)\}.
\] For \(C^1\) \(f\), \(EC(f,K)=\min_{x\in K}m(Df_x)\), where
\(m(Df_x)=\inf_{|v|=1}|Df_xv|\) is the smallest singular value. Hirsch's
\textbf{logarithmic expansion rate} is \[
E(f,K)=\sup_{n\ge0}\frac1n\log EC(f^n,K).
\] This is the maximal asymptotic rate at which \(f\) expands distances
in the fixed Riemannian metric. For a linear expanding map of
\(\mathbb T^n\), \(EC(f)=\min_i|\lambda_i|\) and
\(E(f)=\log\min_i|\lambda_i|\). For a conformal expanding map,
\(EC(f)=\min|f'|\) and \(E(f)=\log\lambda\) for the expansion rate
\(\lambda\). (The definition is quoted in Schreiber 1997.) A note on
kind: \(EC(f,K)\) is metric-bound --- it is a property of the chosen
metric --- while its asymptotic rate \(E(f,K)\) is not (Section 4.3).
Neither is a topological invariant.

\subsection{Schreiber's theorem: the smooth
identification}\label{schreibers-theorem-the-smooth-identification}

Schreiber's theorem is the exact smooth analogue of your
\(\log\lambda_{\exp}=h_{\min}\). Let \(f\) be a \(C^1\) endomorphism of
a Riemannian manifold, \(K\) compact, positively invariant, and \(f\)
immersive at \(K\) (\(\det Df_x\ne0\) for \(x\in K\)). Then \[
E(f,K)=\inf_\mu E(f,\mu),
\] where \(E(f,\mu)\) is the smallest Lyapunov exponent of \((f,\mu)\)
and the infimum is over ergodic \(f\)-invariant measures supported on
\(K\). Equivalently, \[
E(f,K)=\inf_{x\in K}\liminf_{n\to\infty}\frac1n\log m(Df^n_x).
\] Consequences: \(E(f,K)\) is independent of the Riemannian metric
(Lyapunov exponents are), and it is invariant under bi-Lipschitz
conjugacy. For a \(C^{1+\alpha}\) map, \(E(f,K)<-\lambda\) iff there are
exponentially attracted points with rate \(>\lambda\); for a
one-dimensional map, \(E(f,K)<-\lambda\) implies a periodic orbit with
Lyapunov exponent \(<-\lambda\). The paper also answers Hirsch's
questions about aperiodic torus flows (expansion rate zero) and
conjugacies (H\"{o}lder exponents). A note on kind: both sides of
Schreiber's identity are smooth quantities --- they need a
differentiable structure and they are independent of the Riemannian
metric. Your identity \(\log\lambda_{\exp}=h_{\min}\) has the same shape
but lives in the topological column on both sides; the analogy is one of
form, not an instance of a common theorem.

This is the theorem to cite when you want to say ``the maximal expanding
constant is the infimum of the Lyapunov exponents'' in the smooth case.
In your topological setting, \(h_{\min}\) replaces the infimum of the
Lyapunov exponents, and \(\log\lambda_{\exp}=h_{\min}\) is the analogue.

A measure-theoretic analogue is Morales (2025, arXiv:2504.16456): the
expansion exponent of a measurable map is the minimum of the expansion
exponents over Borel probability measures, and a map expands small
distances (in Reddy's sense) iff every Borel probability measure has
positive expansion exponent.

\subsection{\texorpdfstring{Over all metrics: your \(\lambda_{\exp}\)
and its nearest
relatives}{Over all metrics: your \textbackslash lambda\_\{\textbackslash exp\} and its nearest relatives}}\label{over-all-metrics-your-lambda_exp-and-its-nearest-relatives}

Your definition \[
\lambda_{\exp}(f)=\sup\{\lambda>1:\ \exists\text{ compatible }d,\ \epsilon>0\text{ with }d(fx,fy)\ge\lambda d(x,y)\text{ for }d(x,y)<\epsilon\}
\] is a supremum over \textbf{all} compatible metrics. Being a supremum
over all compatible metrics, \(\lambda_{\exp}\) is a topological
conjugacy invariant by construction --- the same device that makes
\(\lambda_{\mathrm{Lip}}\) topological (Section 8.5). I did not find
this exact definition in the literature. The nearest relatives below are
mostly theorems about particular metrics; only their suprema over a
conjugacy-invariant class of metrics are topological. The nearest
relatives are:

\begin{itemize}
\tightlist
\item
  \textbf{Dovbysh's optimal Lyapunov metrics} (2006). For expansive
  homeomorphisms, he constructs compatible Lyapunov metrics with
  prescribed contraction/expansion constants, and for
  A$^{\#}$-homeomorphisms proves that the lower bounds for the constants
  are attained simultaneously. This is the two-sided analogue of
  ``choose the best metric.'' The one-sided/non-invertible case is not
  covered, but the techniques (Lyapunov metrics, Ruelle's Lipschitz
  axiom) are the same circle.
\item
  \textbf{Fujita--Kato--Matsumoto's expanding ratios} (2010). For Ruelle
  expanding maps, they construct metrics with \textbf{exact} expansion
  ratio \(\lambda\) for every \(1<\lambda<s\), and prove the bound
  \(\log\lambda\le h(f)/\dim X\). Their ``strictly expands small
  distances'' is stronger than your inequality expansion, so their set
  of exact ratios is a subset of the \(\lambda\)'s allowed by your
  definition; the suprema may or may not agree. Their Proposition 4.1
  and Remark 4.4 are the cleanest published form of the
  entropy--dimension--expansion sandwich.
\item
  \textbf{Roth--Roth's HausLip constant} (2020). They optimize
  \(\mathrm{HD}_d(X)\log^+\operatorname{Lip}_d(f)\) over metrics and
  prove it equals \(h(f)\) for positively expansive maps. Their
  construction gives a metric with
  \(\operatorname{Lip}_D(f^n)\le16\,\mathrm{SkewLoc}_D(f^n)\), i.e.~a
  near-uniform expansion metric, but it optimizes the product with the
  Hausdorff dimension, not \(\log\lambda\) alone.
\item
  \textbf{Fathi's hyperbolic metric} (1989). For hyperbolic sets, the
  metric that makes the stable/unstable rates uniform; the expanding
  case gives the entropy--dimension inequality.
\item
  \textbf{Schreiber's \(E\)} (1997). In the smooth case, the supremum
  over Riemannian metrics of the asymptotic expansion rate is \(E\),
  because \(E\) is metric-independent. The one-step expansion constant
  \(EC(f,K)\) does depend on the metric, but its asymptotic rate does
  not.
\end{itemize}

\textbf{One inequality between \(\lambda_{\exp}\) and \(E\) is
elementary} (this is your observation, and I have not seen it stated).
Let \(d\) be a Riemannian metric and suppose \(f^n\) is
\((\rho,\epsilon)\)-expanding in \(d\), that is \(EC_d(f^n,K)\ge\rho\).
Put \(\lambda=\rho^{1/n}\) and average the metric over the cycle: \[
d'(x,y)=\sum_{i=0}^{n-1}\lambda^{-i}d(f^ix,f^iy).
\] Then \(d'\) is a compatible metric on a neighbourhood of \(K\), and
for \(d(x,y)\) small, \[
d'(fx,fy)-\lambda\,d'(x,y)=\lambda^{1-n}d(f^nx,f^ny)-\lambda\,d(x,y)\ \ge\ 0,
\] because \(d(f^nx,f^ny)\ge\rho\,d(x,y)=\lambda^n d(x,y)\). So \(f\) is
\((\lambda,\epsilon')\)-expanding in \(d'\), whence
\(\lambda_{\exp}\ge\rho^{1/n}\); taking the supremum over \(n\), \[
\lambda_{\exp}\ \ge\ e^{E(f,K)}.
\] The class inclusion alone gives only
\(\lambda_{\exp}\ge\sup_{d\ \mathrm{Riem}}\lambda_{\mathrm{loc}}(d)\);
the averaging step is what turns ``one step of \(f^n\)'' into ``one step
of \(f\)'' and produces the exponent \(1/n\). The reverse inequality
does not follow this way, and I did not find it stated; by your theorem
it is equivalent to \(h_{\min}\le E(f,K)\).

The draft's theorem \(\log\lambda_{\exp}=h_{\min}\) is, as far as I
found, not in the literature. The closest statement is Schreiber's
\(E=\inf_\mu E(f,\mu)\); the draft's \(h_{\min}\) is the topological
replacement for the infimum of Lyapunov exponents. That is a clean way
to position the theorem: it extends a known smooth identity to expansive
topological systems, with the infimum over measures replaced by a
minimum over pairs of points (or an in-radius growth rate).

\subsection{The entropy--dimension--expansion
sandwich}\label{the-entropydimensionexpansion-sandwich}

A note on kind: every quantity in this subsection is computed in a fixed
metric --- \(D_d(X)\), \(\mathrm{HD}_d(X)\) and \(\log\lambda\) all
depend on \(d\). The inequalities matter because they bound the
topological \(h_{\mathrm{top}}\) by metric data, which is a bridge
between columns rather than an identity inside one.

For a metric \(d\) in which \(f\) expands small distances by at least
\(\lambda\), the lower half of Fujita--Kato--Matsumoto's Proposition 4.1
(which uses only the lower bound) gives \[
D_d(X)\log\lambda\le h_{\mathrm{top}}(f).
\] Since the topological dimension is a lower bound for the box
dimension, \(D_d(X)\ge\dim X\), hence \[
\log\lambda\le\frac{h_{\mathrm{top}}(f)}{\dim X},\qquad
\lambda_{\exp}(f)\le e^{h_{\mathrm{top}}(f)/\dim X}.
\] This is the universal upper bound on the maximal expanding constant.
It is the quantitative form of your \(h_{\min}\le h_{\mathrm{top}}\):
\(\log\lambda_{\exp}=h_{\min}\le h_{\mathrm{top}}\), and the gap
\(h_{\mathrm{top}}/\dim X-h_{\min}\) measures the anisotropy of the
expansion.

Conversely, if \(f\) is Lipschitz with constant \(L\) in the metric
\(d\), the Dai--Zhou--Geng/Misiurewicz inequality gives \[
h_{\mathrm{top}}(f)\le \mathrm{HD}_d(X)\log L.
\] For a metric that is both \(\lambda\)-expanding and \(L\)-Lipschitz,
the two inequalities combine: \[
D_d(X)\log\lambda\;\le\;h_{\mathrm{top}}(f)\;\le\;\mathrm{HD}_d(X)\log L.
\] (Here \(D_d(X)\ge\mathrm{HD}_d(X)\), so the two bounds are
consistent.) If the metric is \textbf{exactly} \(\lambda\)-expanding
(both a lower and an upper bound \(\lambda\)), then
\(h_{\mathrm{top}}(f)=D_d(X)\log\lambda\), which is the ``fractal
metric'' formula of Fujita--Kato--Matsumoto. If it is only
\(\lambda\)-expanding, the lower bound is the relevant one for
\(\lambda_{\exp}\); the upper bound involves the Lipschitz constant,
which may be much larger.

The equality case \(\lambda_{\exp}=e^{h_{\mathrm{top}}/\dim X}\) is what
one expects for \textbf{conformal} expanding maps on a manifold: all
directions expand at the same rate, the Hausdorff dimension of the
adapted metric is the topological dimension, and Bowen's equation
\(h_{\mathrm{top}}=\dim\cdot\log\lambda\) holds. For non-conformal maps,
the best metric can only expand the slowest direction at the slowest
rate, so \(\lambda_{\exp}<e^{h_{\mathrm{top}}/\dim X}\) in general. The
2-cat torus in your draft is the simplest example:
\(\lambda_{\exp}=\sqrt5-1\approx1.236\), while
\(e^{h_{\mathrm{top}}/2}=2\).

\subsection{Examples where the maximal constant is
known}\label{examples-where-the-maximal-constant-is-known}

In every case the invariant being computed is the topological
\(\lambda_{\exp}\); what changes from one example to the next is the
structure used to compute it --- smooth, algebraic, or combinatorial.

\begin{itemize}
\tightlist
\item
  \textbf{Circle doubling / graph maps.} For \(z\mapsto z^d\) on
  \(S^1\), \(\lambda_{\exp}=d=e^{h_{\mathrm{top}}}\); for a graph
  endomorphism with irreducible Markov matrix, the draft's theorem gives
  \(\log\lambda_{\exp}=h_{\min}=h_{\mathrm{top}}=\log\lambda_{\mathrm{Lip}}\),
  so \(\lambda_{\exp}\) is the spectral radius of the Markov matrix.
  This is the dimension-one case where the upper bound is attained. For
  a general \(C^1\) expanding map of \(S^1\) of degree \(d\),
  \(h_{\mathrm{top}}=\log d\) (Misiurewicz--Przytycki), so
  \(\lambda_{\exp}\le d\); equality holds for the constant-slope model,
  while a map such as \(x\mapsto 2x+\epsilon\sin(2\pi x)\) has
  derivative varying between \(2-2\pi\epsilon\) and \(2+2\pi\epsilon\),
  so its maximal expanding constant is generally strictly less than
  \(2\). The entropy--degree formula is Misiurewicz--Przytycki (1977).
\item
  \textbf{Piecewise monotone interval maps.} Parry's constant-slope
  theorem: every transitive piecewise monotone interval map of positive
  entropy \(h\) is conjugate to a map of constant slope \(\pm e^h\)
  (Parry 1966). Hence the maximal expanding constant is
  \(e^{h_{\mathrm{top}}}\) for such a map. Bobok--Roth show more
  generally that for every piecewise monotone interval map and every
  \(C^\infty\) interval map the infimum of Lipschitz constants in the
  conjugacy class equals \(e^{h_{\mathrm{top}}}\), and that the
  inequality can be strict for countably piecewise monotone maps; in the
  topologically mixing Markov case they identify the infimum with the
  exponential of the Salama entropy of a reverse Markov chain.
\item
  \textbf{Linear torus maps.} For \(A\in GL(n,\mathbb Z)\) with all
  eigenvalues \(|\lambda_i|>1\), the best expanding constant is
  \(\min_i|\lambda_i|\) (the slowest eigendirection), so
  \(h_{\min}=\log\min_i|\lambda_i|\), while
  \(h_{\mathrm{top}}=\sum_i\log|\lambda_i|\). The 2-cat matrix
  \(\begin{pmatrix}0&2\\2&2\end{pmatrix}\) has eigenvalues
  \(1\pm\sqrt5\), so \(\lambda_{\exp}=\sqrt5-1\) and
  \(h_{\mathrm{top}}=\log4\). This is my own elementary check, not a
  citation; it matches the draft's example.
\item
  \textbf{Conformal expanding maps.} For a hyperbolic rational map on
  its Julia set, the visual metric gives a metric with
  \(\lambda=e^{h_{\mathrm{top}}/\mathrm{HD}(J)}\) in the conformal
  gauge, so \(\lambda_{\exp}\ge e^{h_{\mathrm{top}}/\mathrm{HD}(J)}\);
  the Ahlfors regular conformal dimension is the infimum of the
  dimensions of the gauge, and the maximal expansion exponent in the
  conformal gauge is \(\log d/\mathrm{ARCdim}\) for a cxc map of degree
  \(d\). The draft's companion paper on hyperbolic quadratic polynomials
  is the place where this is computed; I have not tried to reconcile its
  ``period of the attracting cycle'' statement with the conformal
  formula, since the two may use different normalizations.
\item
  \textbf{Graph virtual endomorphisms.} The draft asserts
  \(h_{\min}=-\log\overline E^\infty[\pi,\phi]\),
  i.e.~\(\lambda_{\exp}=1/\overline E^\infty[\pi,\phi]\), the reciprocal
  of the asymptotic \(\infty\)-conformal energy / stretch factor. This
  is the combinatorial analogue of Parry's constant-slope theorem and is
  the natural place where the maximal expanding constant is exactly
  computable. The stretch-factor background is Bestvina--Handel (1992)
  --- the draft's \texttt{BH92:TrainTracksFreeGp} --- and Thurston's
  1986 norm paper; the draft's own keys for this material are
  \texttt{Thurston16:RubberBands}, \texttt{Thurston20:Characterize} and
  \texttt{PT26:ConfDim}.
\end{itemize}

\subsection{Attainment and optimal
metrics}\label{attainment-and-optimal-metrics}

Your \(\lambda_{\exp}\) is a supremum, and the draft notes that it is
appropriate because \(\lambda\) can be decreased locally near a point.
Whether the supremum is attained is a separate question. The literature
on attainment:

\begin{itemize}
\tightlist
\item
  Dovbysh (2006) proves attainment of the lower bounds for the
  contraction/expansion constants for A$^{\#}$-homeomorphisms, with one
  metric satisfying all the conditions.
\item
  Fujita--Kato--Matsumoto (2010) construct exact expanding metrics for
  every \(\lambda\) below a bound, but they do not claim attainment at
  the supremum; their metrics approach the bound from below.
\item
  Roth--Roth (2020) construct a metric with \(\operatorname{Lip}\) and
  local skew comparable up to a factor 16, which is enough for the
  HausLip equality but not necessarily optimal for \(\lambda\).
\item
  Fathi (1989) constructs a hyperbolic metric for a hyperbolic set,
  again with uniform rates but not necessarily the optimal ones.
\end{itemize}

A natural question (not answered here) is whether \(\lambda_{\exp}\) is
attained for, say, graph maps or expanding maps of the circle. The
examples above suggest it is attained in dimension one, but the general
case is not covered by the papers I found.

\begin{center}\rule{0.5\linewidth}{0.5pt}\end{center}

\section{Other choices in the draft's
appendices}\label{other-choices-in-the-drafts-appendices}

\subsection{Neighborhoods of the diagonal and uniform
spaces}\label{neighborhoods-of-the-diagonal-and-uniform-spaces}

Your Appendix \texttt{alternatives.tex} replaces open covers by open
neighborhoods \(N\) of the diagonal, with the distance defined by
relation composition \(N^k\). This is the \textbf{uniform space}
formulation. The chain distance associated to an entourage is the
standard ``uniformity chain''; the volume versions are the
packing/covering numbers of the uniformity. The entropy side is the
\textbf{uniform entropy} of Bowen: for a uniform space
\((X,\mathcal U)\) and a uniformly continuous self-map, one defines
\(r_n(V,K,\psi)\) as the minimal cardinality of a set \(F\) such that
\(\bigcup_{x\in F}\bigcap_{k<n}\psi^{-k}(V(\psi^k(x)))\) covers \(K\),
and
\(h_{\mathrm{unif}}(\psi)=\sup_V\sup_K\limsup_n\frac1n\log r_n(V,K,\psi)\).
For compact spaces this coincides with topological entropy; for
noncompact spaces it is a uniform invariant. For the uniform entropy see
Kimura (1998) and the survey of Dikranjan--Sanchis--Virili (2012);
Bowen's entropy for uniformly continuous maps is Bowen (1971).

The point for your appendix: the neighborhood-of-the-diagonal language
is not just a convenient reformulation; it is the standard language of
uniform dynamics, and the ``volume'' definitions in that language are
the uniform entropy ones. The distance definitions (relation composition
\(N^k\)) are the uniform chain distance. The two-cover in-radius has a
uniform analogue: the smallest \(r(n)\) such that
\((N^{(n)})^{r(n)}\subseteq N\), which is exactly the uniform version of
your \(r_{\mathcal U}(n)\).

\subsection{\texorpdfstring{Metrics, \(\epsilon\)-nets, and
\(\epsilon\)-chains}{Metrics, \textbackslash epsilon-nets, and \textbackslash epsilon-chains}}\label{metrics-epsilon-nets-and-epsilon-chains}

Your second appendix's metric formulation is the standard Bowen/Dinaburg
setup: \(d_n^f(x,y)=\max_{k<n}d(f^k x,f^k y)\),
\(\epsilon\)-separated/spanning sets, and \(\epsilon\)-chains. The
``volume'' notions are \(\epsilon\)-packing, \(\epsilon\)-covering, and
\(\epsilon\)-nets; the ``distance'' notion is the \(\epsilon\)-chain
length. The chain-based entropy formula of Misiurewicz (1986) is the
relevant prior work if you want to recover the volume entropy from
chains. The chain recurrence literature (Conley theory;
Richeson--Wiseman) uses the same \(\epsilon\)-chains for recurrence and
mixing times, and relates the chain mixing time to entropy.

A caution: the \(\epsilon\)-chain distance \(d(x,y;\epsilon)\) is not
comparable to the metric distance \(d(x,y)\) from above unless the
metric is a length metric. Your appendix notes this; the same caveat
appears in the coarse geometry literature, where \(\epsilon\)-chains
define a large-scale metric only after taking \(\epsilon\to0\) and
waiting for connectivity.

At the large-scale end there is a separate notion of \textbf{coarse
entropy} (W. Geller, M. Misiurewicz, D. Sawicki, ``Coarse entropy of
metric spaces,'' \emph{Geom. Dedicata} \textbf{218} (2024), article 99;
arXiv:2109.15172). It uses \(\delta\)-pseudoorbits and \(R\)-separated
sets with \(\delta,R\to\infty\), is a coarse invariant of the metric
space, and takes only the values \(0\) or \(\infty\). It is a
distance-based entropy, but at large rather than small scales, so it is
not a substitute for your \(h_{\min}\) on a compact space; it is a
natural ``other choice'' if you ever want a noncompact analogue.

\subsection{Connected components of
preimages}\label{connected-components-of-preimages}

Your Remark in \texttt{intro.tex} suggests replacing the iterated cover
\(\mathcal U^{(n)}\) by the connected components of \(f^{-(n+1)}(U)\)
for \(U\in\mathcal U\), at least for locally connected \(X\). This is
exactly the Ha\"{i}ssinsky--Pilgrim cxc setup: for a topologically cxc
map, the preimages of a cover by connected open sets have connected
components that shrink, and the visual metric is built from the
resulting tree of components. The equivalence of the two constructions
(iterated covers vs.~connected components of preimages) is discussed in
HP09, and the local connectivity hypothesis is what makes the components
well-behaved. This is a good ``other choice'' to cite: it is not just a
technical variant, it is the standard way to build visual metrics for
cxc maps.

\subsection{limsup/liminf vs limit}\label{limsupliminf-vs-limit}

Your \(h_{\mathrm{dist}}(x,y;\mathcal U)\) is defined with a limit and
you note that it may not exist; one can replace it by limsup/liminf.
Sun's \(h_L^\pm\) and \(h_E^\pm\) are the standard examples of this
choice; Dai--Jiang's distance entropy also uses critical exponents
rather than a limit. The draft's \(h_{\min}\) uses Fekete's lemma
(supermultiplicativity of \(r_{\mathcal U}(n)\)) to get a limit, which
is the right move for the minimum growth. For \(h_{\mathrm{diam}}\) and
\(h_{\max}\), the LEO hypothesis is what makes the limits exist; this is
analogous to the standard ``submultiplicative if LEO'' arguments in
complex dynamics.

\subsection{Diameter growth, maximum growth, and
LEO}\label{diameter-growth-maximum-growth-and-leo}

Your \(h_{\mathrm{diam}}\) and \(h_{\max}\) are the ``maximal'' end of
the distance spectrum. The smooth analogue is the maximal Lyapunov
exponent / diameter growth of the manifold; the volume analogue is the
Newhouse--Yomdin \(k\)-volume growth. The LEO (locally eventually onto)
hypothesis is standard in one-dimensional dynamics and in the theory of
expanding maps; it ensures that the diameter growth is coarsely
submultiplicative. I did not find a paper that studies exactly your
\(h_{\mathrm{diam}}\) for topological dynamical systems, but the smooth
and conformal cases are covered by the Lyapunov/volume-growth
literature.

\subsection{Polynomial and generalized
growth}\label{polynomial-and-generalized-growth}

If \(h_{\mathrm{top}}=0\), the exponential growth rate is zero and one
can measure polynomial or other growth instead. The relevant notions are
\textbf{polynomial entropy} (M. Marco, ``Polynomial entropies and
integrable Hamiltonian systems,'' \emph{Regul. Chaotic Dyn.} \textbf{18}
(2013), 623--655) and \textbf{generalized entropy} for arbitrary orders
of growth (J. Correa and E. R. Pujals, ``Orders of growth and
generalized entropy,'' arXiv:2003.01257, 2020). The definitions are in
the same spirit as yours: replace the exponential rate by a different
growth scale, and measure the separation of orbits at that scale. This
is a natural ``other choice'' if you want the theory to say something in
the zero-entropy case.

\subsection{Cleaning covers}\label{cleaning-covers}

Your \texttt{intro.tex} asks whether the operation --- take unions of
ascending chains, then delete dominated sets, then restrict to
non-dominated elements --- has a general name. The composite still has
none that I can find, neither for covers nor in lattice theory. Asked
separately, though, \textbf{both halves are standard, and each has a
name}: the chain half is ordinary order theory, the deletion half is the
theory of antichains. Neither name is a name \emph{for covers}, and the
cover literature's own reduction of a cover is a different and stronger
operation. So the answer to ``does this have some general name in
lattice theory?'' is: the vocabulary for both halves exists, the
vocabulary for the composite does not.

\textbf{The chain half: adding upper bounds of ascending chains.} On a
family of subsets, ``adjoin the union of every ascending chain of
members'' is the operation of closing the family under unions of chains.
In that generality it is a well-studied property, with three equivalent
descriptions.

\emph{Order theory.} A poset in which every chain has a least upper
bound is \textbf{chain-complete}, and Markowsky (1976) is the paper that
develops the notion; it proves the non-obvious direction, that a
chain-complete poset is a \textbf{dcpo} --- every directed subset has a
least upper bound --- so ``closed under unions of chains'' and ``closed
under unions of directed subfamilies'' are one condition, not two (for
families of sets this is Mayer-Kalkschmidt--Steiner 1964). The same
paper gives the universal form of your operation: the \textbf{chain
completion} of a poset is the free chain-complete poset on it, and
Markowsky says explicitly that it is not the MacNeille completion. Two
remarks that matter for your setting. First, nothing needs completing:
\(\mathcal O(X)\) is already a complete lattice, so the operation is the
plain one of closing a subfamily under the chain suprema that already
exist. Second, a chain-complete \emph{lattice} is already a complete
lattice (Markowsky's introduction), so for a lattice the chain
hypothesis is not a weak one: it already yields all suprema.

\emph{Universal algebra.} For families closed under intersections the
dictionary is exact. Nation's Theorem 3.1: for a closure operator
\(\Gamma\), the following are equivalent --- \(\Gamma\) is
\textbf{algebraic} (finitary); the union of any up-directed family of
\(\Gamma\)-closed sets is \(\Gamma\)-closed; the union of any
\emph{chain} of \(\Gamma\)-closed sets is \(\Gamma\)-closed. So the
chain step is exactly what makes a closure system algebraic, and the
adjective for a family closed under directed unions in that literature
is \textbf{inductive} (Ern\'{e}'s term).

\emph{Convexity.} A family closed under arbitrary intersections and
under directed unions is exactly a \textbf{convexity} in van de Vel's
sense, and the theory of convex structures (van de Vel 1993) is where
``closed under directed unions'' is an axiom rather than a lemma.

\emph{For covers specifically} I found no name and no use of this
closure. Two standard facts sit next to it. First, in your setting the
chain step is nearly vacuous at the top: \(X\) is compact if and only if
the top element of the frame \(\mathcal O(X)\) is a compact element
(Picado--Pultr 2012), i.e.~if and only if every directed family of open
sets whose union is \(X\) already contains \(X\); so if
\(X\notin\mathcal U\), no chain-union of members is \(X\). Second, the
classical Zorn argument for covers does take a union along an ascending
chain, but the chain is a chain of \emph{subcollections}, not of
members: Arens and Dugundji (1950) call a subcollection redundant if the
remaining members still cover, note that the union of an ascending chain
of redundant subcollections is redundant --- this is where
point-finiteness enters --- and conclude with Zorn's lemma that every
open cover of a metacompact space has a minimal cover refining it.

\textbf{The deletion half: deleting dominated sets.} Keeping only the
members contained in no other member produces what combinatorics calls a
\textbf{Sperner family}, an \textbf{antichain}, or a \textbf{clutter}.
The map \(\max\) (``take the maximal members'') needs no hypothesis to
be \emph{defined}, and it is always monotone, idempotent and
contractive, i.e.~an interior operator. What needs a hypothesis is
\textbf{uniqueness}, and you were right to ask for one: the
correspondence the counting literature states --- down-set
\(\leftrightarrow\) antichain of its maximal members --- is a bijection
only under a chain condition, and finiteness is the special case usually
meant.

\begin{itemize}
\tightlist
\item
  \emph{What holds without hypotheses.} For \emph{finitely many}
  antichains, ``unite, then delete the dominated members'' is the join
  in the order ``every member of \(\alpha\) lies in some member of
  \(\beta\)'': \(\alpha\vee\beta=\max(\alpha\cup\beta)\), written
  \(\lceil\cdot\rceil\) in the antichain literature (De
  Wulf--Doyen--Henzinger--Raskin 2006). No finiteness of the ambient
  poset is needed here, because a chain in a union of \(k\) antichains
  has at most \(k\) elements, so nothing can be lost.
\item
  \emph{What needs the hypothesis.} Every member of the family must lie
  below a \emph{maximal} member of the family. Finiteness, and more
  generally the ascending chain condition, gives this. In general the
  hypothesis is exactly the one you named --- every ascending chain in
  the family has an upper bound in the family, which is Zorn's
  hypothesis. Without it the maximal set carries no information: on
  \([0,1]\), the families \(\{(0,1-1/n)\}_{n\ge2}\) and
  \(\{(0,1-1/n)\}_{n\ge2}\cup\{(0,1/2)\}\) both have empty maximal set,
  and the family \(\{\{0,\dots,n\}\}_{n\ge0}\) of finite subsets of
  \(\mathbb N\) has no maximal member at all. For infinitely many
  antichains the same failure kills the join: the antichains
  \(\{\{0,\dots,n\}\}\) have no upper bound at all among the antichains
  of the finite subsets of \(\mathbb N\), since no finite set contains
  every \(\{0,\dots,n\}\), and therefore no least one.
\item
  \emph{Kleitman's conditions are exactly the finite case.} His objects
  are the subsets of a finite \(n\)-element set \(S_n\), and his ideal
  is a down-set (\(B\in I\) and \(A\subseteq B\) implies \(A\in I\)).
  Then ``an ideal can be uniquely corresponded to an antichain (its
  maximal members)'' holds --- in a finite poset every member of a
  down-set lies below a maximal member \emph{of that down-set} --- and
  the count is the Dedekind number. When \(P\) is the Boolean lattice on
  \(n\) atoms, that lattice of antichains is the \textbf{free
  distributive lattice on \(n\) generators}, whose elements are the
  monotone Boolean functions (Birkhoff 1967, p.~63; Kleitman 1969). So
  Kleitman's sentence is your step \emph{with a hypothesis supplied},
  and the hypothesis he uses is finiteness, not a chain condition. That
  is why the sentence does not transfer to covers of a compact space as
  it stands: a cover of \([0,1]\) can contain an infinite ascending
  chain with no top, and then its antichain has lost a member.
\end{itemize}

For covers, the standard reduction is stronger, and this is the part
worth getting right in the paper. The literature's notion is not ``no
member contains another'' but \textbf{minimal}, or \textbf{irreducible}:
no proper subfamily covers, equivalently no member is contained in the
union of the \emph{others} (Arens--Dugundji 1950; a space is called
irreducible when every open cover has an irreducible open refinement).
Deleting dominated members is strictly weaker. On a three-point space
the cover \(\{\{1,2\},\{2,3\},\{1,3\}\}\) has no member containing
another, yet any two of its members cover. Irreducibility implies the
Sperner condition and not conversely, so your deletion step is the
\emph{antichain} (Sperner) reduction of a cover, not the
\emph{irreducible} (minimal-subcover) reduction: on an irreducible cover
the antichain reduction does nothing, and on an antichain cover it need
not be irreducible. Nor is your deletion step ``deleting redundant
members'' in the sense of Arens--Dugundji: a member can be redundant
without being contained in any single other member, which is the same
three-point example.

\textbf{What this means for the draft.} (1) If you want to name the
composite rather than each half, ``the antichain of maximal members of
the chain-closure'' is accurate but not a term of art; there is nothing
to cite, and one sentence saying so is the right amount. (2) The two
steps are not independent conveniences: the chain step is what makes the
deletion step legitimate, by supplying precisely the chain condition
above. From \(\mathcal U\) alone the maximal members need not cover ---
on \([0,1]\), \(\mathcal U=\{(0,1-1/n)\}_{n\ge2}\cup\{[0,1/2],\{1\}\}\)
has maximal members \([0,1/2]\) and \(\{1\}\), which miss \((1/2,1)\)
--- whereas with the chain-closure taken in the idempotent sense, Zorn's
lemma puts every member below a maximal chain-union, so the maximal
chain-unions do cover. (That argument is my reading of your definition,
not a theorem quoted from a source.) (3) One caution about the
definition as written: condition (1) quantifies over \(\mathcal U\), not
over the chain-closure, and then the output need not be an antichain. On
\([0,1]\) take
\(\mathcal U=\{(0,1/2)\}\cup\{(1/j,1)\}_{j\ge2}\cup\{\{0\},\{1/2\},\{1\}\}\);
it is a cover, and both \((0,1/2)\) and \((0,1)=\bigcup_j(1/j,1)\)
satisfy your two conditions, nested. Quantifying condition (1) over the
chain-closure makes the output exactly the set of maximal chain-unions,
which is an antichain and, with the closure taken in the idempotent
sense, still covers \(X\). (4) A last simplification you can state
without proof: a \emph{finite} cover is already chain-complete, so for
finite covers your cleaning is just the deletion of dominated sets.

\begin{center}\rule{0.5\linewidth}{0.5pt}\end{center}

\section{What I did not find}\label{what-i-did-not-find}

To be explicit about the boundary of this search:

\begin{itemize}
\tightlist
\item
  I did \textbf{not} find a paper that defines your \(\lambda_{\exp}\)
  as the supremum over \textbf{all compatible metrics} of the local
  expansion rate. The closest are Dovbysh's optimal Lyapunov metrics
  (two-sided, expansive homeomorphisms), Fujita--Kato--Matsumoto's exact
  expanding ratios (Ruelle expanding maps), and Schreiber's
  metric-independent asymptotic rate (smooth maps).
\item
  I did \textbf{not} find a paper that proves your identity
  \(\log\lambda_{\exp}=h_{\min}\) in the topological setting. The
  closest is Schreiber's \(E=\inf_\mu E(f,\mu)\) for \(C^1\) maps. This
  is the natural ``prior work'' boundary for the theorem.
\item
  I did \textbf{not} find a paper that studies your pairwise growth
  spectrum \(h_{\mathrm{dist}}(x,y)\) as a function of \((x,y)\) in the
  cover-chain setting. The smooth analogue is the Lyapunov spectrum /
  multifractal analysis; the cover analogue may be new.
\item
  I did \textbf{not} find a paper that uses the in-radius
  \(r(\mathcal U;\mathcal V)\) as a primary invariant. Sun's Lebesgue
  number is the closest, and the comparison
  \(r\asymp\delta/\operatorname{mesh}\) is my own derivation, not a
  quoted theorem.
\item
  I did \textbf{not} find a paper that defines your
  \(h_{\mathrm{diam}}\) or \(h_{\max}\) in the cover-chain sense, nor
  one that asserts or denies \(h_{\max}=\log\lambda_{\mathrm{Lip}}\).
  The closest are Kifer's pointwise maximum/minimum characteristic
  exponents (metric-dependent, no supremum over metrics) and the
  adapted-metric theorems behind the smooth analogue; see Section 8.
\item
  I did \textbf{not} find a paper that studies the infimum of the
  Lipschitz constant over \emph{length} metrics, as opposed to over all
  compatible metrics (where it is trivially \(\le1\)), over a conjugacy
  class of Euclidean metrics (Bobok--Roth), or over metrics in a product
  with Hausdorff dimension (Roth--Roth).
\item
  I did \textbf{not} find any paper that treats the composite operation
  of your \texttt{intro.tex} --- close a cover under unions of ascending
  chains, then keep the maximal members --- as an object with a name, or
  that uses it. The two halves are separately standard and separately
  named (Section 5.7); the composite is not.
\item
  I did \textbf{not} read every paper in full. The statements above are
  from the abstracts, introductions, and the specific theorems I cite; I
  have not verified every proof. The Fujita--Kato--Matsumoto inequality
  is from their Proposition 4.1/Remark 4.4; the Sun inequalities are
  from the two Sun papers; the Schreiber theorem is from Schreiber's
  paper. The comparisons labelled ``my own'' are just that.
\end{itemize}

\begin{center}\rule{0.5\linewidth}{0.5pt}\end{center}

\section{Follow-up: Julia sets, polynomial Julia sets, hyperbolic
quadratic
polynomials}\label{follow-up-julia-sets-polynomial-julia-sets-hyperbolic-quadratic-polynomials}

This section answers the follow-up question specifically. It
deliberately leaves out the conformal-dimension story
(Ha\"{i}ssinsky--Pilgrim, ARCdim, the Pilgrim--Thurston energy paper),
which is already known. I have separated work from your own group
(Braga--Sieg, Section 7.1) from independent prior work (Hruska, Section 7.2;
Gelfert--Przytycki--Rams, Section 7.3). A note on kind: everything in this
section is smooth --- metrics on \(J\), derivatives, Lyapunov exponents.
The one exception is the Braga--Sieg \(\lambda(f)\): it is a supremum
over metrics on \(J\), so as an invariant of \((J,f)\) it is
topological, even though every metric used to compute it is conformal.

\subsection{Braga--Sieg: the exact sup-over-metrics definition on a
Julia set (work from your
group)}\label{bragasieg-the-exact-sup-over-metrics-definition-on-a-julia-set-work-from-your-group}

The closest match to your \(\lambda_{\exp}\) on a Julia set is
\textbf{Vitor Braga and David Sieg, ``Expansivity on Julia sets''} (2023
Indiana University REU report, 20 pp.,
\url{https://math.indiana.edu/documents/reu/2023/Braga-Sieg.pdf};
presented at YMC 2024 as David R. Sieg, ``Expansivity on Julia Sets''; I
did not find a journal version). They define, for a rational map
\(f:\widehat{\mathbb C}\to\widehat{\mathbb C}\) with Julia set \(J\):

\begin{itemize}
\tightlist
\item
  \(f\) is \((\lambda,\epsilon)\)-expansive on \(J\) with respect to a
  metric \(d\) if for all \(z_1,z_2\in J\) with \(d(z_1,z_2)<\epsilon\),
  one has \(d(fz_1,fz_2)\ge \lambda d(z_1,z_2)\);
\item
  \(f\) is \(\lambda\)-expansive on \(J\) if this holds for some
  \(\epsilon>0\);
\item
  the \textbf{expansivity} of \(f\) on \(J\) is \[
  \lambda(f)=\sup\{\lambda:\ \exists\text{ a metric }d\text{ on }J\text{ with }f\ (\lambda,\epsilon)\text{-expansive on }J\}.
  \]
\end{itemize}

This is exactly your \(\lambda_{\exp}(f)\) restricted from \(X\) to the
Julia set \(J\), with the same ``sup over compatible metrics'' and the
same local inequality. It is the only place I found where this exact
definition is written down for rational maps.

They relate \(\lambda(f)\) to the asymptotic \(\infty\)-energy
\(E^\infty\) of the associated virtual endomorphism / spine. Their two
key lemmas (stated in the report, with proofs deferred to a forthcoming
research paper) are:

\begin{itemize}
\tightlist
\item
  \textbf{Lemma 14:} if \(f^n\) has an \(f^n\)-invariant graph
  \(G_0\subset G_n\) with induced map \(g:G_0\to G_0\) and Markov
  partition matrix \(A\), then \[
  \lambda(f)\ \le\ \bigl(\lambda_n(g)\bigr)^{1/n},
  \] the \(n\)-th root of the Perron--Frobenius eigenvalue.
\item
  \textbf{Lemma 15:} if \(E^\infty[\pi,\phi]<1\), then \[
  E^\infty[\pi,\phi]\ \ge\ \frac{1}{\lambda(f)}.
  \] With the invariant graph of Lemma 14,
  \(E^\infty[\pi,\phi]\ge 1/(\lambda_n(g))^{1/n}\).
\end{itemize}

Thus invariant graphs give upper bounds for \(\lambda(f)\), and
\(E^\infty\) gives lower bounds; when the two match,
\(\lambda(f)=1/E^\infty\). This is the same relation as the draft's
asserted \(h_{\min}=-\log E^\infty\) in the comb-models section, with
\(\log\lambda(f)=h_{\min}\).

\textbf{Quadratic examples.} They compute \(E^\infty\) for several
quadratic census maps of Brezin--Byrne--Levy--Pilgrim--Plummer:

\begin{itemize}
\tightlist
\item
  \textbf{Quadratic map \#7:} \(E^\infty=2^{-1/2}\), so the matching
  bounds give \(\lambda(f)=\sqrt2\). (The report lists its critical
  portrait as type C(4,4); I have not decoded the notation, so I will
  not label the cycle period.)
\item
  \textbf{Quadratic map \#3.2} (period-3 critical cycle, the rabbit
  side): \(E^\infty=2^{-1/3}\), so the matching bounds give
  \(\lambda(f)=2^{1/3}\). This is exactly the ``\(\lambda_{\exp}\)
  records the period of the attracting cycle'' phenomenon in a specific
  period-3 case.
\item
  \textbf{Quadratic map \#6.1} (a more complicated critical portrait):
  \(E^\infty\approx0.739\), so the matching bounds give
  \(\lambda(f)\approx1.353\) (if the displayed bounds are tight; the
  report's exact algebraic equation for this case is at the end of Section 4).
\item
  \textbf{Cubic map \#51:} the report only gets
  \(0.553<E^\infty<0.563\); this is an example where the period alone
  does not determine the answer. (Lemma 15 turns the upper bound into
  \(\lambda(f)>1.78\); an upper bound for \(\lambda(f)\) would need a
  separate invariant-graph estimate.)
\end{itemize}

Caveats: this is an REU report, not a refereed paper; Lemmas 14--15 are
announced with proofs deferred; and the examples are quadratic
\emph{census} maps, which include polynomial cases (the period-2 and
period-3 critical-cycle cases above) but are stated for quadratic
rational maps. Since Braga and Sieg were working under your supervision,
this is your group's own work, not independent prior work. I record it
here as the exact internal match to your definition; the independent
comparison is Hruska in Section 7.2.

\subsection{Hruska: explicit expanding metrics and box-expansion
constants}\label{hruska-explicit-expanding-metrics-and-box-expansion-constants}

The other direct line is \textbf{Suzanne Lynch Hruska, ``Constructing an
expanding metric for dynamical systems in one complex variable,''}
\emph{Nonlinearity} \textbf{18} (2005), no. 1, 81--100;
arXiv:math/0406003; and her Cornell thesis, \textbf{``On the numerical
construction of hyperbolic structures for complex dynamical systems''}
(August 2002), available at
\url{https://www.math.stonybrook.edu/theses/thesis02-2/part1.pdf}.

The paper's setup is the standard one: a polynomial
\(f:\mathbb C\to\mathbb C\) of degree \(>1\) is hyperbolic if there is a
Riemannian metric \(\mu\) on a neighborhood \(B\) of \(J\) and an
expansion constant \(\lambda>1\) such that
\(|D_zf(v)|_\mu\ge\lambda|v|_\mu\) for every \(z\in B\) and
\(v\in T_z\mathbb C\). Hruska introduces a \textbf{box metric}: a
piecewise constant conformal metric \(\Phi\) on a finite box cover
\(\{B_k\}\) of \(J\), with ``handicaps'' \(\phi_k>0\). The map is
\textbf{box-expansive} by \(L>1\) if, for every edge \(B_k\to B_j\) of
the box graph, every \(z\in B_k\) and every \(v\in T_z\mathbb C\), \[
\phi_j\,|D_zf(v)|\ \ge\ L\,\phi_k\,|v|.
\] Theorem 1.4 proves that box expansion by \(L>1\) implies a smooth
Riemannian metric on a neighborhood of \(J\) with expansion at least
\(L\). So every computed \(L\) is a \textbf{lower bound} for the maximal
expanding constant of \(f\) on \(J\).

The key combinatorial result is Proposition 1.7: for a box graph
\(\Gamma\), let \(\underline L\) be the minimum, over all simple cycles
\(B_0\to\cdots\to B_{n-1}\to B_0\), of the average cycle multiplier \[
\bigl(\lambda_0\cdots\lambda_{n-1}\bigr)^{1/n},\qquad
\lambda_k=\min_{z\in B_k}|f'(z)|.
\] Then a box metric with box expansion \(L\) exists \textbf{iff}
\(L\le\underline L\). So \(\underline L\) is the ideal box-expansion
constant on that finite graph; as the grid refines, the minimum average
cycle multiplier should converge to the infimum of Lyapunov exponents
over periodic orbits, i.e.~the smooth analogue of \(e^{h_{\min}}\). I
have not found a paper that states this convergence for Julia sets
explicitly; it is the symbolic-image / Morse-spectrum convergence of
Osipenko (Hruska cites Osipenko {[}17,18{]} as the related computational
method).

\textbf{Quadratic examples.} Hruska computes explicit box-expansion
constants (all lower bounds for \(\lambda_{\exp}\) on \(J\)):

{\def\LTcaptype{none} % do not increment counter
\begin{longtable}[]{@{}
  >{\raggedright\arraybackslash}p{(\linewidth - 6\tabcolsep) * \real{0.2500}}
  >{\raggedright\arraybackslash}p{(\linewidth - 6\tabcolsep) * \real{0.2500}}
  >{\raggedright\arraybackslash}p{(\linewidth - 6\tabcolsep) * \real{0.2500}}
  >{\raggedright\arraybackslash}p{(\linewidth - 6\tabcolsep) * \real{0.2500}}@{}}
\toprule\noalign{}
\begin{minipage}[b]{\linewidth}\raggedright
Polynomial
\end{minipage} & \begin{minipage}[b]{\linewidth}\raggedright
Parameter
\end{minipage} & \begin{minipage}[b]{\linewidth}\raggedright
Box expansion constant \(L\)
\end{minipage} & \begin{minipage}[b]{\linewidth}\raggedright
Source
\end{minipage} \\
\midrule\noalign{}
\endhead
\bottomrule\noalign{}
\endlastfoot
basilica & \(z^2-1\) & \(1.14067\) & paper and thesis \\
aeroplane & \(z^2-1.755\) & \(1.05069\) & paper and thesis \\
Cantor Julia set & \(z^2+0.35\) & \(1.00778\) & paper \\
rabbit & \(z^2-0.123+0.745i\) & \(1.00508\) (subdivided \(2^8\) grid) &
thesis \\
double basilica & \(z^2-1.31\) & \(1.06758\) & thesis \\
five-eared rabbit & \(z^2-0.504+0.563i\) & \(1.0264\) & thesis \\
six-eared rabbit & \(z^2+0.388+0.22i\) & not completed (memory limit) &
thesis \\
\end{longtable}
}

The published paper's examples are the basilica, aeroplane, and Cantor
set; the thesis has the rabbit, double basilica, and five/six-eared
rabbits. The computed constants are lower bounds, not exact
\(\lambda(f)\); they depend on the box model. The thesis's rabbit value
\(1.00508\) is much smaller than the \(2^{1/3}\approx1.260\) suggested
by the Braga--Sieg period-3 computation, which is consistent with it
being a coarse lower bound rather than the supremum.

\subsection{The Lyapunov-spectrum
line}\label{the-lyapunov-spectrum-line}

Outside the polynomial setting, the same ``sup over visual metrics''
question appears for expanding Thurston maps. Bonk--Meyer define visual
metrics on the sphere and their expansion factors; Yin proves that the
supremum of the expansion factors over visual metrics is the
combinatorial expansion factor, that it is attained for Latt\`{e}s maps,
and that it need not be attained for Latt\`{e}s-type maps (Q. Yin,
``Latt\`{e}s maps and combinatorial expansion,'' \emph{Ergodic Theory
Dynam. Systems} \textbf{36} (2016), 1307--1342; Bonk--Meyer,
\emph{Expanding Thurston Maps}, AMS Math. Surveys Monogr. 225, 2017).
This is the same ``maximize expansion over a natural class of metrics''
idea, but over visual metrics rather than all compatible metrics.

For a hyperbolic rational map, the range of Lyapunov exponents on \(J\)
is a natural ``smooth'' version of the distance-growth spectrum.
\textbf{K. Gelfert, F. Przytycki, M. Rams, ``On the Lyapunov spectrum
for rational maps,''} \emph{Math. Ann.} \textbf{348} (2010), no. 4,
965--1004; arXiv:0809.3363, prove that for a rational map with no
critical points in \(J\), the Lyapunov spectrum is described by the
pressure function and has a range \([\alpha^-,\alpha^+]\). The lower
endpoint \(\alpha^-\) is the infimum of the Lyapunov exponents on \(J\),
so \(e^{\alpha^-}\) is the maximal expansion rate in the smooth
conformal category (the analogue of \(\lambda_{\exp}\) for a hyperbolic
polynomial). The paper applies to hyperbolic quadratic polynomials but
does not compute the endpoint for specific parameters.

Measure-specific computations:

\begin{itemize}
\tightlist
\item
  \textbf{S. Sridharan and A. R. Tiwari, ``The dependence of Lyapunov
  exponents of polynomials on its coefficients,''} arXiv:1906.03864
  (2019): computes Lyapunov exponents of hyperbolic quadratic and cubic
  polynomials with respect to various Bernoulli measures, using
  thermodynamic formalism. This is not the infimum over all measures,
  but it is a computational line for expansion rates on Julia sets.
\item
  \textbf{``Lyapunov exponents of polynomials with respect to certain
  weighted Lyubich's measures,''} arXiv:1908.07711: another
  measure-specific computation for hyperbolic quadratic polynomials.
\item
  \textbf{The lower Lyapunov exponent at the critical value}
  (Levin--Przytycki--Shen, arXiv:1304.6178) is a different quantity: it
  is negative iff the unicritical polynomial has an attracting cycle. It
  is not the Julia-set expansion constant.
\end{itemize}

\subsection{What I did not find}\label{what-i-did-not-find-1}

\begin{itemize}
\tightlist
\item
  I did \textbf{not} find a published journal paper that defines
  \(\lambda(f)\) as the supremum over metrics on \(J\) and computes it
  for hyperbolic quadratic polynomials. The only exact-definition match
  is the Braga--Sieg REU report (unpublished), which is your group's own
  work (supervised by you), not independent prior work.
\item
  I did \textbf{not} find a paper that states the general formula
  \(\lambda(f)=2^{1/p}\) for a hyperbolic quadratic polynomial with an
  attracting cycle of period \(p\). The Braga--Sieg values \(2^{-1/2}\)
  and \(2^{-1/3}\) are consistent with it, and their
  ``\(2^{-1/p}\le E^\infty\)'' lower bound for a critical cycle of type
  \(C(2,p)\) is the mechanism, but the general statement is not in the
  sources I found. This is likely the gap the draft's companion paper
  fills.
\item
  Hruska's constants are \textbf{lower bounds}, not exact
  \(\lambda(f)\); the exact values would require the ideal box constant
  in the limit, which the paper does not compute.
\item
  I did not find any independent (non-REU, non-Thurston-circle) paper
  that uses the sup-over-metrics expansion constant on Julia sets.
  Hruska's box-expansion constant is the closest independent
  computational notion, and it is a lower bound rather than the
  supremum.
\end{itemize}

\begin{center}\rule{0.5\linewidth}{0.5pt}\end{center}

\section{Follow-up: the other end --- diameter growth, maximum growth,
and the minimal Lipschitz
constant}\label{follow-up-the-other-end-diameter-growth-maximum-growth-and-the-minimal-lipschitz-constant}

Your memory is right: the maximum-end material is in
\texttt{clippings.tex} and never reached \texttt{min-growth.tex}. As you
put it, it was formerly part of the paper and was cut when the planned
scope changed, which is why it now reads as an unintegrated sketch ---
it carries a TODO list rather than an argument. It sits directly after
your note ``Start on maximum growth stuff, assuming LEO. This turns out
not to be sufficient for submultiplicativity, and it's not the focus of
this paper'', so the definitions there are what the comparison below is
against.

The short answer: there is a well-developed literature for the
\textbf{pointwise} maximum exponent in a metric space (Kifer 1983 and
its descendants), a small one for \textbf{infima of Lipschitz constants}
(Bobok--Roth 2019, Roth--Roth 2020), and a large one for the
\textbf{volume}-growth version (Newhouse, Yomdin). I found no paper that
defines \(h_{\mathrm{diam}}\) or \(h_{\max}\) in the cover-chain sense,
and none that asserts or denies \(h_{\max}=\log\lambda_{\mathrm{Lip}}\).

\subsection{\texorpdfstring{What the draft has (quoted from
\texttt{clippings.tex})}{What the draft has (quoted from clippings.tex)}}\label{what-the-draft-has-quoted-from-clippings.tex}

\begin{itemize}
\tightlist
\item
  \textbf{Diameter growth.}
  ``\(h_{\mathrm{diam}}(f;\mathcal U)=\lim_{n\to\infty}\frac{\operatorname{diam}(X;\mathcal U_n)}{n}\)
  (if it exists)'', with the note ``(Something about refinement)''. The
  displayed formula is missing the logarithm: if
  \(\operatorname{diam}(X;\mathcal U_n)\) is the diameter in the chain
  distance \(d(\cdot,\cdot;\mathcal U_n)\), that diameter grows like
  \(\lambda^n\) in every expanding example, so the quotient without
  \(\log\) diverges and its limit is not the growth rate; it reads as a
  typo for \(\frac1n\log\operatorname{diam}(X;\mathcal U_n)\).
\item
  \textbf{Maximum growth.} Recorded as a sketch rather than a
  definition: ``Define \(h_{\max}\) as supremum over covers
  \(\mathcal U\) of diameter of \(\mathcal U_n\cap U\), for
  \(U\in\mathcal U_n\). Only finite if \(X\) is locally connected.''
\item
  \textbf{Lipschitz stretch.} ``\(\lambda_{\mathrm{Lip}}\) is the
  infimum, over all length metrics on \(X\), of the best Lipschitz
  constant for \(f\)'', together with your remark that the length-metric
  restriction is necessary because snowflaking \(d\mapsto d^\epsilon\)
  (\(\epsilon<1\)) decreases the Lipschitz constant.
\item
  \textbf{The one inequality.}
  ``\(h_{\max}\le\log\lambda_{\mathrm{Lip}}\) when defined.''
\item
  \textbf{The chain.}
  \(h_{\min}\le h_{\mathrm{diam}}\le h_{\max}\le h_{\mathrm{top}}\) for
  \(f\) locally eventually onto, expansive, \(X\) compact locally
  connected, all inequalities strict in examples; and the figure-eight
  example (\(z\mapsto z^3\) on one loop, \(z\mapsto z^2\) on the other)
  in which \(\lambda_{\mathrm{Lip}}\) and \(e^{h_{\max}}\) drop from
  \(3\) to \(2\) when passing to the \(z\mapsto z^2\) subsystem.
\end{itemize}

\subsection{Where the diameter analysis
went}\label{where-the-diameter-analysis-went}

The comparison of \(h_{\mathrm{diam}}\) with \(h_{\max}\) --- global
versus localised growth, what the two give for a linear expanding
endomorphism, and the 2Cat example --- is analysis rather than prior
work, so it is a separate note,
\texttt{reports/diameter-growth-2cat.md}. Its conclusion, after several
corrections: \(h_{\mathrm{diam}}=\log2\) for the 2Cat with the ball
cover (the map factors as \(g\,h\) with \(g=\operatorname{diag}(2,2)\)
and \(h\in GL_2(\mathbb Z)\); the chain-adjacency coincide exactly under
\(h^{n-1}\)), while \(h_{\max}=\log2\tau\) is the local eigen-expansion
--- so the four-value chain
\(h_{\min}<h_{\mathrm{diam}}<h_{\max}<h_{\mathrm{top}}\) of
\texttt{clippings.tex} is intact and strict. What belongs in this memo
is the literature, which is Section 8.3--8.7 below.

\subsection{Kifer 1983: the closest published
match}\label{kifer-1983-the-closest-published-match}

\textbf{Yuri Kifer, ``Characteristic exponents of dynamical systems in
metric spaces,''} \emph{Ergodic Theory Dynam. Systems} \textbf{3}
(1983), no. 1, 119--127; DOI 10.1017/S0143385700001838. For a
one-parameter group \(S^t\) of homeomorphisms of a separable metric
space \(X\) without isolated points (discrete or continuous time), he
sets \[
B_x(\delta,T)=\{y\in X\setminus\{x\}:\ d(S^tx,S^ty)<\delta\ \text{for all } t\in[0,T]\},
\] \[
A_\delta(x,t)=\sup_{y\in B_x(\delta,t)}\frac{d(S^tx,S^ty)}{d(x,y)},\qquad
a_\delta(x,t)=\inf_{y\in B_x(\delta,t)}\frac{d(S^tx,S^ty)}{d(x,y)},
\] and then \[
\Lambda^\pm_\delta(x)=\limsup_{t\to\pm\infty}\frac{1}{|t|}\ln A_\delta(x,t),\qquad
\lambda^\pm_\delta(x)=\limsup_{t\to\pm\infty}\frac{1}{|t|}\ln a_\delta(x,t).
\] Under an integrability hypothesis on a Borel invariant probability
measure \(\mu\) (satisfied whenever \(S^t\) is uniformly Lipschitz for
\(|t|\le1\)), Theorem 1 gives that \(\mu\)-almost everywhere the limsups
are limits and that \(\Lambda^+_\delta(x)=-\lambda^-_\delta(x)\),
\(\lambda^+_\delta(x)=-\Lambda^-_\delta(x)\); and if the hypothesis
holds for all small \(\delta\), the limits \[
\Lambda^\pm(x)=\lim_{\delta\to0}\Lambda^\pm_\delta(x),\qquad \lambda^\pm(x)=\lim_{\delta\to0}\lambda^\pm_\delta(x)
\] exist. Theorem 1 also gives the invariance of
\(\Lambda^\pm,\lambda^\pm\) under \(S^t\) and their constancy
\(\mu\)-a.e. when \(\mu\) is ergodic. Theorem 2 does the same for an
invariant set \(Q\), with
\(B_Q(\delta,T)=\{y\in X\setminus Q: d(S^ty,Q)<\delta\ \forall t\in[0,T]\}\)
and \(d(y,Q)\) in the denominator --- the exact analogue of localizing
at a point, one level up. Theorem 3 is a stability criterion:
\(\Lambda^+(x)<0\) plus forward Lyapunov stability implies asymptotic
stability, while \(\Lambda^+(x)>0\) excludes forward stability. Theorem
4 and its Corollary identify the quantities in the smooth case: \[
\Lambda^+(x)=\lim_{t\to\infty}\frac1t\ln\|DS^t_x\|,\qquad
\lambda^+(x)=\lim_{t\to\infty}\frac1t\ln\inf_{\xi\in T_x,\ \|\xi\|=1}\|DS^t_x\xi\|,
\] \(\mu\)-a.e., so that ``\(\Lambda^+(x)\) and \(\lambda^+(x)\) are the
maximal and the minimal characteristic exponents at \(x\)''. Remark 1 is
the reason this is the right comparison for you: \(\Lambda^+(x)>0\) and
\(\lambda^+(x)<0\) is ``a very weak kind of hyperbolicity''. Remark 2
matters for a one-sided setting: the positive-time half of Theorem 1
holds when \(S^t\) is only a semigroup.

\textbf{How it matches the draft.} \(A_\delta(x,t)\) is exactly ``sup
over \(y\) of the distance ratio at time \(t\), among the \(y\) whose
orbit shadows that of \(x\) to accuracy \(\delta\)'' --- the metric
counterpart of \(\sup_y\) of the pairwise growth in
\(\mathcal U\)-chains, with ``shadows to accuracy \(\delta\)'' replacing
``lies in the same \(\mathcal U_t\)-atom''. So \(\sup_x\Lambda^+(x)\) is
the metric-space form of \(h_{\max}\), and \(h_{\mathrm{diam}}\) is what
one gets by moving the \(\sup\) over \(x\) outside (\emph{and} dropping
the shadowing restriction), which is a different and cruder object. The
comparison between the chain and metric formulations is my reading, not
a theorem in Kifer's paper --- it is the same kind of comparison as
\(r(\mathcal U;\mathcal V)\asymp\delta(\mathcal U)/\operatorname{mesh}(\mathcal V)\)
in Section 3, and it should be stated with the same hedge.

\textbf{What Kifer does not do.} He never takes a supremum over
compatible metrics, so his \(\Lambda^\pm\) are the analogues of your
\emph{fixed-cover} quantities, not of the topological invariants. And
the quantity is genuinely metric-dependent: see Section 8.4. Kind: metric-bound
--- \(\Lambda^\pm\) are unchanged by a bi-Lipschitz change of the metric
and changed by a topological conjugacy; the topological content is the
sign.

\subsection{The pointwise-exponent line: metric dependence, and what
survives
it}\label{the-pointwise-exponent-line-metric-dependence-and-what-survives-it}

\textbf{C. A. Morales, P. Thieullen, H. Villavicencio, ``Lyapunov
exponents on metric spaces,''} \emph{Bull. Austral. Math. Soc.}
\textbf{97} (2018), 153--162; DOI 10.1017/S0004972717000703. They
replace Kifer's two limits by the pointwise Lipschitz constant: for
\(f:X\to X\) and \(x\) not isolated, \[
f'(x)=\limsup_{y\to x}\frac{d(fx,fy)}{d(x,y)},\qquad
\chi_f(x)=\limsup_{n\to\infty}\frac1n\log (f^n)'(x).
\] Theorem 4.2: on a Riemannian manifold with the Riemannian distance,
\(\chi_f(x)=\Lambda_f(x)=\limsup_n\frac1n\log\|Df^n_x\|\) --- the
classical upper exponent. Theorem 3.2: at a Lyapunov-stable point \(x\),
\(\chi_f(x)\le\Delta^+(f,x)\) (their rendering of Kifer's quantity);
Remark 3.3 says they do not know whether the inequality can be strict.
Theorem 5.1: a periodic point with \(\chi_f(x)<0\) is asymptotically
stable. Corollary 5.4: for an expansive homeomorphism of a compact
metric space without isolated points there is a \emph{compatible} metric
in which \(\inf_x\max\{\chi_f(x),\chi_{f^{-1}}(x)\}>0\) --- Fathi's
metric again.

Two corrections to the way this line is usually summarized, both of
which matter for your definition:

\begin{itemize}
\tightlist
\item
  Their invariance is \textbf{bi-Lipschitz}, not topological. Theorem
  4.4(3) says \(\chi_{hfh^{-1}}=\chi_f\circ h^{-1}\) for their
  ``diffeomorphisms'', which are homeomorphisms \(h\) with
  \(\|h'\|_\infty<\infty\) and \(\|(h^{-1})'\|_\infty<\infty\). Example
  2.6 states the boundary explicitly: bi-Lipschitz equivalent metrics
  give the same \(\chi_f\), and ``this may be false if the metrics are
  only compatible''. So \(\chi_f\) is an invariant of the Lipschitz
  structure, not of the topological conjugacy class.
\item
  The metric dependence is not a technicality. \textbf{M. J. Pacifico,
  J. L. Vieitez, ``Lyapunov exponents for expansive homeomorphisms,''}
  \emph{Proc. Edinburgh Math. Soc.} (2020), DOI
  10.1017/S0013091519000579; arXiv:1704.05284, work in Kifer's framework
  and prove that on a compact Peano space an expansive homeomorphism has
  \(\Lambda^+(x)\ge\gamma>0\) and \(\lambda^+(x)\le-\gamma\) for some
  \(\gamma>0\) and \emph{every} \(x\), in a suitable metric --- while
  giving an example where the original metric does not see this. What
  expansiveness pins down is the \textbf{sign}, in some metric; not the
  value, and not in a fixed metric.
\end{itemize}

This is the precise sense in which a \emph{topological} maximum-end
invariant must either be built from covers (your \(h_{\max}\)) or be an
infimum/supremum over metrics (your \(\lambda_{\mathrm{Lip}}\)). The
pointwise exponents in the literature are neither, which is why they
read as ``almost'' your objects.

\subsection{\texorpdfstring{\(\lambda_{\mathrm{Lip}}\): why the
length-metric restriction is doing real work, and the closest published
relatives}{\textbackslash lambda\_\{\textbackslash mathrm\{Lip\}\}: why the length-metric restriction is doing real work, and the closest published relatives}}\label{lambda_mathrmlip-why-the-length-metric-restriction-is-doing-real-work-and-the-closest-published-relatives}

\textbf{An elementary remark on the definition.} Without a restriction
on the metric the definition collapses. For any continuous self-map of a
compact metric space, any compatible \(d\) and any \(\epsilon>0\), \[
d_\epsilon(x,y)=\sum_{m\ge0}(1+\epsilon)^{-m}d(f^mx,f^my)
\] is a compatible metric and
\(\operatorname{Lip}_{d_\epsilon}(f)\le 1+\epsilon\). (Check:
\(d_\epsilon(fx,fy)=\sum_{m\ge0}(1+\epsilon)^{-m}d(f^{m+1}x,f^{m+1}y)=(1+\epsilon)\sum_{k\ge1}(1+\epsilon)^{-k}d(f^kx,f^ky)\le(1+\epsilon)d_\epsilon(x,y)\).)
So the infimum over \emph{all} compatible metrics is at most \(1\) ---
for any map whatsoever --- and it is \(0\) for a constant map. The
length-metric restriction is therefore not a technical convenience: it
is the only thing preventing \(\lambda_{\mathrm{Lip}}\) from being
trivial, and your snowflake remark is a second face of the same fact.
For a homeomorphism the same computation with weights \(2^{-|n|}\) gives
an invariant metric with Lipschitz constant \(1\), again non-length in
general. (This paragraph is my own computation, not a quotation.)

\textbf{Jozef Bobok, Samuel Roth, ``The infimum of Lipschitz constants
in the conjugacy class of an interval map,''} \emph{Proc. Amer. Math.
Soc.} \textbf{147} (2019), 255--269; arXiv:1612.06095. They study \[
\Lambda(f)=\inf\{\lambda:\ \exists\ \text{a homeomorphism }\psi:[0,1]\to[0,1]\text{ with }\psi f\psi^{-1}\ \lambda\text{-Lipschitz}\},
\] an infimum over the pullbacks of the Euclidean metric by
homeomorphisms --- not over length metrics, and not over all compatible
metrics. Kind: \(\Lambda(f)\) is a \textbf{topological} invariant,
because the infimum runs over a conjugacy-invariant class of metrics;
unlike a Lipschitz constant in a fixed metric, it sees only the
conjugacy class. Your \(\lambda_{\mathrm{Lip}}\) is topological for the
same reason, with length metrics in place of pullbacks of the Euclidean
metric. Results:

\begin{itemize}
\tightlist
\item
  \(\Lambda(f)\ge e^{h(f)}\) for positive entropy (their Section 2; the general
  lower bound).
\item
  \(\Lambda(f)\le\nu(f)\) (Theorem 3.4), where
  \(\nu(f)=\limsup_n(\operatorname{Var}f^n)^{1/n}\) is the growth rate
  of the \textbf{total variation} of the iterates --- not the number of
  iterated preimages. (The number of preimages enters a different
  theorem; see the next point. I flag this because it is easy to
  conflate the two, and one summary of this paper does.)
\item
  \(\Lambda(f)=e^{h(f)}\) for piecewise monotone maps (Corollary 3.5)
  and for \(C^\infty\) maps (Corollary 3.6).
\item
  For countably piecewise monotone, topologically mixing, Markov maps,
  \[
  \log\Lambda(f)=h_{\mathrm{RevSal}}(\Sigma)=\limsup_{n\to\infty}\frac1n\log\#f^{-n}(\{x\})
  \] (Theorem 4.17), the Salama entropy of the reverse Markov chain,
  equal to the growth rate of the number of iterated preimages of a
  nearly arbitrary point. Example 4.19 shows this can be
  \textbf{strictly larger} than \(h(f)\).
\item
  If \(f\) is locally eventually onto --- their ``leo'' is your LEO for
  interval maps --- then \(\Lambda(f)=e^{h(f)}\) again (Theorem 4.15).
\end{itemize}

So the closest published relative of \(\lambda_{\mathrm{Lip}}\) always
satisfies \(\ge e^{h}\), equals \(e^{h}\) exactly when the combinatorics
is tame (piecewise monotone, \(C^\infty\), or leo), and exceeds
\(e^{h}\) only for maps with unbounded combinatorics. That is a
suggestive parallel to your chain
\(h_{\min}\le h_{\mathrm{diam}}\le h_{\max}\le h_{\mathrm{top}}\), and
it is only a parallel: the two infima are taken over different classes
of metrics, so there is no implication either way.

\textbf{J. Milnor, \emph{Dynamics: one-dimensional}, Section 13}
(\url{https://www.math.stonybrook.edu/~jack/DYNOTES/dn13.pdf}).
The one place I found where the ``sup over metrics'' move is stated as a
programme, and it is stated at the \emph{minimum} end: ``Here the
constant \(c\) depends on the particular choice of Riemannian metric.
However the expansion constant \(k\) is independent of the metric.
\ldots{} (To obtain a differentiable conjugacy invariant, we could take
the supremum of all such expansion constants.)'' It is a remark, not a
theorem, but it is the smooth counterpart of your \(\lambda_{\exp}\) and
it points at the same design question at the other end.

\textbf{S. Roth, Z. Roth, ``Inequalities for entropy, Hausdorff
dimension, and Lipschitz constants,''} \emph{Studia Math.} \textbf{250}
(2020), 253--264; arXiv:1807.01196, already in Section 3, is the other
infimum-over-metrics construction:
\(\mathrm{HausLip}(X,f)=\inf_d \mathrm{HD}_d(X)\log^+\operatorname{Lip}_d(f)\),
a product rather than a bare Lipschitz constant, which is a topological
conjugacy invariant and equals \(h(f)\) for positively expansive maps
and linear torus maps. Its method --- construct a metric in which the
Lipschitz constant is controlled uniformly over iterates --- is the
technology any proof of an inequality between \(h_{\max}\) and
\(\log\lambda_{\mathrm{Lip}}\) would use.

\subsection{The volume end: where the two ends meet in one
inequality}\label{the-volume-end-where-the-two-ends-meet-in-one-inequality}

\textbf{S. E. Newhouse, ``Entropy and volume,''} \emph{Ergodic Theory
Dynam. Systems} \textbf{8} (1988), 283--299, and \textbf{Y. Yomdin,
``Volume growth and entropy,''} \emph{Israel J. Math.} \textbf{57}
(1987), 285--300. Write \(V_{\ell,k}(f)\) for the supremum of
\(\limsup_n\frac1n\log\omega(f^n\circ\sigma)\) over \(C^k\) maps
\(\sigma\) of the unit \(\ell\)-cube into \(M\), where \(\omega\) counts
\(\ell\)-volume with multiplicity, and \(V_\ell(f)\) for the resulting
\(\ell\)-volume growth rate; and set \[
R(f)=\lim_{n\to\infty}\frac1n\log\max_{x\in M}\|Df^n(x)\|.
\] Then (Newhouse) \(h(f)\le V_{\ell}(f)\) for \(C^{1+\epsilon}\) maps,
and (Yomdin) \(V_{\ell,k}(f)\le h(f)+\frac{2\ell}{k}R(f)\) for \(C^k\)
maps, so for \(C^\infty\) maps the two together give \(h(f)=V_\ell(f)\).
Everything here is smooth: \(V_{\ell,k}\) and \(R\) are built from the
derivative and the Riemannian volume, and are independent of the choice
of Riemannian metric. (I quote the form in L.-S. Young's survey
``Entropy in dynamical systems'', in \emph{Entropy}, Princeton Univ.
Press, 2003, which states Newhouse's and Yomdin's theorems together.)

Two things are worth extracting. First, \(R(f)\) \textbf{is} a
diameter-growth quantity: \(\frac1n\log\) of the maximum over \(x\) of
the operator norm at time \(n\), i.e.~a maximum over \(x\) taken
\emph{at each \(n\)} and then a limit --- a \(\lim_n\sup_x\), not a
\(\sup_x\lim_n\). That is the same order of quantifiers as your
\(h_{\max}\) (a supremum over the atoms of \(\mathcal U_n\) at each
\(n\), then a limit), and it is not the same as \(\sup_x\Lambda^+(x)\);
the two differ exactly when the limsup is not uniform in \(x\). Second,
Yomdin's inequality is the one place in the literature where entropy and
a maximum-end quantity appear in the same inequality:
\(V_{\ell,k}\le h+\frac{2\ell}{k}R\). It says the volume growth is
bounded by entropy plus the maximum distance growth, and for
\(C^\infty\) maps the volume end collapses onto entropy. That is the
closest published analogue I found of your chain, though it runs in the
other direction (volume above entropy, bounded by entropy plus the max
end).

\subsection{What is missing from the
literature}\label{what-is-missing-from-the-literature}

\begin{itemize}
\tightlist
\item
  \textbf{No cover-chain definitions.} I did not find a paper defining
  \(h_{\mathrm{diam}}\) or \(h_{\max}\) as you do. The phrase ``diameter
  growth'' is unusable as a search term: it returns Abresch--Gromoll and
  Riemannian geometry, Schreier-graph diameter growth in group theory,
  and manifold topology. ``Maximal Lyapunov exponent'', ``maximal
  characteristic exponent'' and ``diameter of iterates'' are the working
  terms, and they lead to Kifer.
\item
  \textbf{No comparison theorem.} Nothing I found states
  \(h_{\max}=\log\lambda_{\mathrm{Lip}}\), in either direction beyond
  your Proposition. Kifer and Morales--Thieullen--Villavicencio give
  pointwise, metric-dependent exponents; Bobok--Roth give the infimum in
  a conjugacy class of interval maps; Roth--Roth give an
  infimum-over-metrics product. None of them optimizes over length
  metrics, and none compares the result to a cover-based maximum growth.
\item
  \textbf{No comparison theorem, and no counterexample either.} Nothing
  found states \(h_{\max}=\log\lambda_{\mathrm{Lip}}\) in either
  direction. What the missing inequality would be, and where a
  counterexample would have to live, are set out in
  \texttt{reports/diameter-growth-2cat.md}, Section 7.
\end{itemize}

\begin{center}\rule{0.5\linewidth}{0.5pt}\end{center}

\section{Bibliography}\label{bibliography}

Every work cited in this memo, collected once and grouped by thread.
Within each theme the order is alphabetical by first author. A work that
belongs to two threads is listed under the primary one and named in the
other theme's lead sentence. Passing mentions in the text that carry no
year --- the multifractal papers of Barreira--Pesin--Schmeling,
Pesin--Weiss and Iommi--Jordan in Section 3.7, Kozlovski's theorem on expansion
entropy, Osipenko's computational method, Conley's chain recurrence ---
are not indexed here.

\textbf{A. Entropy from distances: the classical layer}

\begin{itemize}
\tightlist
\item
  \textbf{R. L. Adler, A. G. Konheim, M. H. McAndrew}, ``Topological
  entropy,'' \emph{Trans. Amer. Math. Soc.} \textbf{114} (1965),
  309--319.
\item
  \textbf{R. Bowen}, ``Entropy for group endomorphisms and homogeneous
  spaces,'' \emph{Trans. Amer. Math. Soc.} \textbf{153} (1971),
  401--414.
\item
  \textbf{R. Bowen}, ``Topological entropy for noncompact sets,''
  \emph{Trans. Amer. Math. Soc.} \textbf{184} (1973), 125--136.
\item
  \textbf{E. I. Dinaburg}, ``On the relations among various entropy
  characteristics of dynamical systems,'' \emph{Math. USSR Izv.}
  \textbf{5} (1971), 337--378.
\item
  \textbf{W. H. Gottschalk, G. A. Hedlund}, \emph{Topological Dynamics},
  AMS Colloquium Publications \textbf{36}, 1955.
\item
  \textbf{P. Walters}, \emph{An Introduction to Ergodic Theory},
  Graduate Texts in Mathematics \textbf{79}, Springer, 1982.
\end{itemize}

\textbf{B. Lebesgue numbers, expansive constants, in-radii}

\begin{itemize}
\tightlist
\item
  \textbf{F. Przytycki, M. Urba\'{n}ski}, \emph{Conformal Fractals:
  Ergodic Theory Methods}, London Math. Soc. Lecture Note Series
  \textbf{371}, Cambridge Univ. Press, 2010 (Ch. 4, including Section 4.6,
  ``Expansive maps are expanding in some metric'').
\item
  \textbf{W. L. Reddy}, ``Expanding maps on compact metric spaces,''
  \emph{Topology Appl.} \textbf{13} (1982), 327--334.
\item
  \textbf{M. Shub}, ``Endomorphisms of compact differentiable
  manifolds,'' \emph{Amer. J. Math.} \textbf{91} (1969), 175--199.
\item
  \textbf{P. Sun}, ``Exponential decay of Lebesgue numbers,''
  \emph{Discrete Contin. Dyn. Syst.} \textbf{32} (2012), 3773--3785;
  arXiv:1002.0549.
\item
  \textbf{P. Sun}, ``Exponential decay of expansive constants,''
  arXiv:1101.0851 (2011).
\end{itemize}

\textbf{C. Chains, uniform spaces, and other growth scales}

\begin{itemize}
\tightlist
\item
  \textbf{J. Correa, E. R. Pujals}, ``Orders of growth and generalized
  entropy,'' arXiv:2003.01257 (2020).
\item
  \textbf{X. Dai, Y. Jiang}, ``Distance entropy of dynamical systems on
  noncompact-phase spaces,'' \emph{Discrete Contin. Dyn. Syst.}
  \textbf{20} (2008), 313--333; DOI 10.3934/dcds.2008.20.313.
\item
  \textbf{D. Dikranjan, M. Sanchis, S. Virili}, ``New and old facts
  about entropy in uniform spaces and topological groups,''
  \emph{Topology Appl.} \textbf{159} (2012), no. 7; DOI
  10.1016/j.topol.2011.05.046.
\item
  \textbf{W. Geller, M. Misiurewicz, D. Sawicki}, ``Coarse entropy of
  metric spaces,'' \emph{Geom. Dedicata} \textbf{218} (2024), art. 99;
  arXiv:2109.15172.
\item
  \textbf{T. Kimura}, ``Completion theorem for uniform entropy,''
  \emph{Comment. Math. Univ. Carolin.} \textbf{39} (1998).
\item
  \textbf{M. Marco}, ``Polynomial entropies and integrable Hamiltonian
  systems,'' \emph{Regul. Chaotic Dyn.} \textbf{18} (2013), 623--655.
\item
  \textbf{M. Misiurewicz}, ``Remark on the definition of topological
  entropy,'' in \emph{Dynamical Systems and Partial Differential
  Equations} (Caracas, 1984), Univ. Sim\'{o}n Bol\'{i}var, Caracas,
  1986, 65--67.
\item
  \textbf{D. Richeson, J. Wiseman}, ``Chain recurrence rates and
  topological entropy,'' \emph{Topology Appl.} \textbf{156} (2008),
  251--261; arXiv:0801.1635.
\end{itemize}

\textbf{D. The maximal expanding constant}

\begin{itemize}
\tightlist
\item
  \textbf{M. Bestvina, M. Handel}, ``Train tracks and automorphisms of
  free groups,'' \emph{Ann. of Math.} \textbf{135} (1992), 1--51.
\item
  \textbf{X. Dai, Z. Zhou, X. Geng}, ``Some relations between
  Hausdorff-dimensions and entropies,'' \emph{Sci. China Ser. A}
  \textbf{41} (1998), 1068--1075.
\item
  \textbf{S. A. Dovbysh}, ``Optimal Lyapunov metrics of expansive
  homeomorphisms,'' \emph{Izv. Math.} \textbf{70} (2006), 883--929; DOI
  10.1070/IM2006v070n05ABEH002332.
\item
  \textbf{A. Fathi}, ``Expansiveness, hyperbolicity and Hausdorff
  dimension,'' \emph{Comm. Math. Phys.} \textbf{126} (1989), 249--262.
\item
  \textbf{C. Fujita, H. Kato, M. Matsumoto}, ``Fractal metrics of Ruelle
  expanding maps and expanding ratios,'' \emph{Topology Appl.}
  \textbf{157} (2010), 615--628; DOI 10.1016/j.topol.2009.11.002.
\item
  \textbf{M. W. Hirsch}, ``Asymptotic phase, shadowing, and
  reaction-diffusion systems,'' in \emph{Differential Equations,
  Dynamical Systems, and Control Science}, Lecture Notes in Pure and
  Appl. Math. \textbf{152}, Dekker, 1994, 87--99.
\item
  \textbf{C. A. Morales}, ``A measure-theoretic expansion exponent,''
  arXiv:2504.16456 (2025).
\item
  \textbf{M. Misiurewicz}, ``On Bowen's definition of topological
  entropy,'' \emph{Discrete Contin. Dyn. Syst.} \textbf{10} (2004),
  827--833.
\item
  \textbf{M. Misiurewicz, F. Przytycki}, ``Topological entropy and
  degree of smooth mappings,'' \emph{Bull. Acad. Polon. Sci. S\'{e}r.
  Sci. Math. Astronom. Phys.} \textbf{25} (1977), 573--574.
\item
  \textbf{M. Misiurewicz, W. Szlenk}, ``Entropy of piecewise monotone
  mappings,'' \emph{Studia Math.} \textbf{67} (1980), 45--63.
\item
  \textbf{W. Parry}, ``Symbolic dynamics and transformations of the unit
  interval,'' \emph{Trans. Amer. Math. Soc.} \textbf{122} (1966),
  368--378.
\item
  \textbf{S. Roth, Z. Roth}, ``Inequalities for entropy, Hausdorff
  dimension, and Lipschitz constants,'' \emph{Studia Math.} \textbf{250}
  (2020), 253--264; arXiv:1807.01196.
\item
  \textbf{S. J. Schreiber}, ``Expansion rates and Lyapunov exponents,''
  \emph{Discrete Contin. Dyn. Syst.} \textbf{3} (1997), 433--438; DOI
  10.3934/dcds.1997.3.433.
\item
  \textbf{W. P. Thurston}, ``A norm for the homology of 3-manifolds,''
  \emph{Mem. Amer. Math. Soc.} \textbf{59} (1986), no. 339.
\end{itemize}

The definitions this theme rests on --- Reddy (1982),
Przytycki--Urba\'{n}ski (2010), Shub (1969) --- are listed under B. The
closest relative of \(\lambda_{\mathrm{Lip}}\), Bobok--Roth (2019), is
under E, and Fathi (1989) is also used in G.

\textbf{E. The maximum end: diameter growth, maximum growth, and the
minimal Lipschitz constant}

\begin{itemize}
\tightlist
\item
  \textbf{J. Bobok, S. Roth}, ``The infimum of Lipschitz constants in
  the conjugacy class of an interval map,'' \emph{Proc. Amer. Math.
  Soc.} \textbf{147} (2019), 255--269; arXiv:1612.06095.
\item
  \textbf{Y. Kifer}, ``Characteristic exponents of dynamical systems in
  metric spaces,'' \emph{Ergodic Theory Dynam. Systems} \textbf{3}
  (1983), no. 1, 119--127; DOI 10.1017/S0143385700001838.
\item
  \textbf{J. Milnor}, \emph{Dynamics: one-dimensional}, Stony Brook
  lecture notes, Section 13 (``Expanding maps''),
  \url{https://www.math.stonybrook.edu/~jack/DYNOTES/dn13.pdf}.
\item
  \textbf{C. A. Morales, P. Thieullen, H. Villavicencio}, ``Lyapunov
  exponents on metric spaces,'' \emph{Bull. Austral. Math. Soc.}
  \textbf{97} (2018), 153--162; DOI 10.1017/S0004972717000703.
\item
  \textbf{M. J. Pacifico, J. L. Vieitez}, ``Lyapunov exponents for
  expansive homeomorphisms,'' \emph{Proc. Edinburgh Math. Soc.} (2020);
  DOI 10.1017/S0013091519000579; arXiv:1704.05284.
\end{itemize}

The other infimum-over-metrics construction, Roth--Roth (2020), is
listed under D.

\textbf{F. Volume growth and Lyapunov exponents}

\begin{itemize}
\tightlist
\item
  \textbf{B. R. Hunt, E. Ott}, ``Defining chaos,'' \emph{Chaos}
  \textbf{25} (2015), 097618; arXiv:1501.07896.
\item
  \textbf{S. E. Newhouse}, ``Entropy and volume,'' \emph{Ergodic Theory
  Dynam. Systems} \textbf{8} (1988), 283--299.
\item
  \textbf{V. I. Oseledets}, ``A multiplicative ergodic theorem: Lyapunov
  characteristic numbers for dynamical systems,'' \emph{Trans. Moscow
  Math. Soc.} \textbf{19} (1968), 197--221.
\item
  \textbf{Ya. B. Pesin}, ``Lyapunov characteristic exponents and smooth
  ergodic theory,'' \emph{Russian Math. Surveys} \textbf{32} (1977),
  55--114.
\item
  \textbf{R. Sacksteder, M. Shub}, ``Entropy of a differentiable map,''
  \emph{Adv. Math.} \textbf{28} (1978), 181--185.
\item
  \textbf{Y. Yomdin}, ``Volume growth and entropy,'' \emph{Israel J.
  Math.} \textbf{57} (1987), 285--300.
\item
  \textbf{L.-S. Young}, ``Entropy in dynamical systems,'' in
  \emph{Entropy}, Princeton Univ. Press, 2003.
\end{itemize}

\textbf{G. Visual metrics, conformal dimension, and energies}

\begin{itemize}
\tightlist
\item
  \textbf{M. Bonk, O. Schramm}, ``Embeddings of Gromov hyperbolic
  spaces,'' \emph{Geom. Funct. Anal.} \textbf{10} (2000), 266--306.
\item
  \textbf{R. Bowen}, ``Hausdorff dimension of quasicircles,''
  \emph{Publ. Math. Inst. Hautes \'{E}tudes Sci.} \textbf{50} (1979),
  11--25.
\item
  \textbf{S. Buyalo, V. Schroeder}, \emph{Elements of Asymptotic
  Geometry}, EMS Monographs in Mathematics, European Mathematical
  Society, 2007.
\item
  \textbf{W. J. Floyd}, ``Group completions and limit sets of Kleinian
  groups,'' \emph{Invent. Math.} \textbf{57} (1980), 205--218.
\item
  \textbf{M. Gromov}, ``Hyperbolic groups,'' in \emph{Essays in Group
  Theory}, MSRI Publications \textbf{8}, Springer, 1987, 75--263.
\item
  \textbf{P. Ha\"{i}ssinsky, K. M. Pilgrim}, ``Coarse expanding
  conformal dynamics,'' \emph{Ast\'{e}risque} \textbf{325} (2009).
\item
  \textbf{W. Parry, M. Pollicott}, ``Zeta functions and the periodic
  orbit structure of hyperbolic dynamics,'' \emph{Ast\'{e}risque}
  \textbf{187--188} (1990).
\item
  \textbf{K. M. Pilgrim, D. P. Thurston}, ``Ahlfors-regular conformal
  dimension and energies of graph maps,'' arXiv:2112.09041; \emph{Groups
  Geom. Dyn.}, DOI 10.4171/ggd/914.
\item
  \textbf{D. Ruelle}, ``Repellers for real analytic maps,''
  \emph{Ergodic Theory Dynam. Systems} \textbf{2} (1982), 99--107.
\end{itemize}

\textbf{H. Julia sets and expanding metrics in one complex variable}

\begin{itemize}
\tightlist
\item
  \textbf{M. Bonk, D. Meyer}, \emph{Expanding Thurston Maps}, AMS
  Mathematical Surveys and Monographs \textbf{225}, 2017.
\item
  \textbf{V. Braga, D. Sieg}, ``Expansivity on Julia sets,'' 2023
  Indiana University REU report,
  \url{https://math.indiana.edu/documents/reu/2023/Braga-Sieg.pdf}.
\item
  \textbf{K. Gelfert, F. Przytycki, M. Rams}, ``On the Lyapunov spectrum
  for rational maps,'' \emph{Math. Ann.} \textbf{348} (2010), 965--1004;
  arXiv:0809.3363.
\item
  \textbf{S. L. Hruska}, ``Constructing an expanding metric for
  dynamical systems in one complex variable,'' \emph{Nonlinearity}
  \textbf{18} (2005), 81--100; arXiv:math/0406003; and the thesis
  \emph{On the numerical construction of hyperbolic structures for
  complex dynamical systems}, Cornell University, 2002,
  \url{https://www.math.stonybrook.edu/theses/thesis02-2/part1.pdf}.
\item
  \textbf{G. Levin, F. Przytycki, W. Shen}, ``The Lyapunov exponent of
  holomorphic maps,'' arXiv:1304.6178.
\item
  \textbf{S. Sridharan, A. R. Tiwari}, ``The dependence of Lyapunov
  exponents of polynomials on its coefficients,'' arXiv:1906.03864
  (2019).
\item
  \textbf{S. Sridharan, A. R. Tiwari}, ``Lyapunov exponents of
  polynomials with respect to certain weighted Lyubich's measures,''
  arXiv:1908.07711 (2019).
\item
  \textbf{Q. Yin}, ``Latt\`{e}s maps and combinatorial expansion,''
  \emph{Ergodic Theory Dynam. Systems} \textbf{36} (2016), 1307--1342.
\end{itemize}

\textbf{I. Covers, refinement, and uniform spaces}

\begin{itemize}
\tightlist
\item
  \textbf{R. Engelking}, \emph{General Topology}, revised and completed
  edition, Sigma Series in Pure Mathematics \textbf{6}, Heldermann
  Verlag, Berlin, 1989.
\item
  \textbf{J. R. Isbell}, \emph{Uniform Spaces}, Mathematical Surveys
  \textbf{12}, American Mathematical Society, Providence, RI, 1964.
\item
  \textbf{J. L. Kelley}, \emph{General Topology}, D. Van Nostrand,
  Princeton, NJ, 1955; reprinted as Graduate Texts in Mathematics
  \textbf{27}, Springer, 1975.
\item
  \textbf{S. Willard}, \emph{General Topology}, Addison-Wesley, Reading,
  MA, 1970.
\end{itemize}

\textbf{J. Metric entropy and the partition order}

\begin{itemize}
\tightlist
\item
  \textbf{A. N. Kolmogorov}, ``A new metric invariant of transitive
  dynamical systems and automorphisms of Lebesgue spaces,'' \emph{Dokl.
  Akad. Nauk SSSR} \textbf{119} (1958), 861--864.
\item
  \textbf{V. A. Rokhlin}, ``Lectures on the entropy theory of
  measure-preserving transformations,'' \emph{Russian Math. Surveys}
  \textbf{22} (1967), no. 5, 1--52; Russian original in \emph{Uspekhi
  Mat. Nauk} \textbf{22} (1967), no. 5, 3--56.
\item
  \textbf{Ya. G. Sinai}, ``On the concept of entropy of a dynamical
  system,'' \emph{Dokl. Akad. Nauk SSSR} \textbf{124} (1959), 768--771.
\end{itemize}

\textbf{K. Covers as a lattice: chains, antichains, and the two halves
of ``cleaning''}

\begin{itemize}
\tightlist
\item
  \textbf{R. Arens, J. Dugundji}, ``Remark on the concept of
  compactness,'' \emph{Portugaliae Math.} \textbf{9} (1950), 141--143.
\item
  \textbf{G. Birkhoff}, \emph{Lattice Theory}, 3rd ed., American
  Mathematical Society Colloquium Publications \textbf{25}, Providence,
  RI, 1967 (free distributive lattices and antichains, p.~63).
\item
  \textbf{M. De Wulf, L. Doyen, T. A. Henzinger, J.-F. Raskin},
  ``Antichains: a new algorithm for checking universality of finite
  automata,'' in \emph{Computer Aided Verification (CAV 2006)}, Lecture
  Notes in Comput. Sci. \textbf{4144}, Springer, 2006, 17--30.
\item
  \textbf{D. Kleitman}, ``On Dedekind's problem: the number of monotone
  Boolean functions,'' \emph{Proc. Amer. Math. Soc.} \textbf{21} (1969),
  677--682.
\item
  \textbf{G. Markowsky}, ``Chain-complete posets and directed sets with
  applications,'' \emph{Algebra Universalis} \textbf{6} (1976), 53--68.
\item
  \textbf{J. Mayer-Kalkschmidt, E. Steiner}, ``Some theorems in set
  theory and applications in the ideal theory of partially ordered
  sets,'' \emph{Duke Math. J.} \textbf{31} (1964), 287--289.
\item
  \textbf{J. B. Nation}, \emph{Notes on Lattice Theory}, Ch. 3 (Theorem
  3.1),
  \url{https://math.hawaii.edu/~jb/math618/os3uh.pdf}.
\item
  \textbf{J. Picado, A. Pultr}, \emph{Frames and Locales: Topology
  without Points}, Frontiers in Mathematics, Birkh\"{a}user/Springer,
  Basel, 2012.
\item
  \textbf{M. L. J. van de Vel}, \emph{Theory of Convex Structures},
  North-Holland Mathematical Library \textbf{50}, North-Holland,
  Amsterdam, 1993.
\end{itemize}

\end{document}
